\documentclass[11pt,a4paper]{article}
\usepackage[margin=2.5cm]{geometry}
\usepackage{amsmath,amssymb,amsthm}
\usepackage{enumitem}
\usepackage[hidelinks]{hyperref}
\usepackage[switch]{lineno}
\setlength{\emergencystretch}{2em}
\newtheorem{theorem}{Theorem}[section]
\newtheorem{lemma}[theorem]{Lemma}
\newtheorem{proposition}[theorem]{Proposition}
\newtheorem{corollary}[theorem]{Corollary}
\newtheorem{definition}[theorem]{Definition}
\theoremstyle{remark}
\newtheorem{remark}[theorem]{Remark}
\theoremstyle{plain}
\newcommand{\D}{\mathcal{D}}
\newcommand{\WitnessSet}{\mathsf{W}}
\newcommand{\cf}{\mathrm{cf}}
\newcommand{\lin}{\mathrm{lin}}
\newcommand{\Ccf}[1]{\mathcal C_{#1}^{\cf}}
\newcommand{\Clin}[1]{\mathcal C_{#1}^{\lin}}
\newcommand{\pref}{\mathrm{pref}}
\newcommand{\suf}{\mathrm{suf}}
\title{Distributional Learning of Context-Free Languages under Fixed Finite-Monoid Typing}
\author{Takayuki Kuriyama\\
Independent Researcher\\
\texttt{growup.kuriyama@gmail.com}}
\date{}
\begin{document}
\linenumbers
\maketitle

\begin{abstract}
We study positive-data learning of context-free languages under a fixed
homomorphism $h:\Sigma^*\to M$ into a finite monoid. A language is
$\sim_h$-substitutable if nonempty factors of the same $h$-type that share
one context have identical sets of contexts. For each fixed $h$, we define
a set-driven reconstruction operator that recovers every context-free
$\sim_h$-substitutable target from finitely many positive witnesses, and a
separate conservative sequential learner that identifies every nonempty
target in the limit; reconstruction and updates take polynomial time.
Our characteristic-data bound depends on shortest yields in a target grammar
refined by $h$-values, which can be exponentially larger than ordinary
grammar thickness in some presentations. The classical characterization
of locally trivial finite semigroups gives an exact criterion: $h$ is
determined on nonempty factors by a fixed prefix--suffix window exactly when
$h(\Sigma^+)$ is locally trivial. In this regime ordinary grammar thickness
suffices, and the resulting union of language classes coincides with
Yoshinaka's fixed-window hierarchy. For linear targets we
obtain characteristic data polynomial in linear grammar size under every
fixed $h$. We also exhibit nonregular linear and nonlinear context-free
targets outside every fixed-window class, while the one-bracket Dyck language
lies outside every fixed-typing class.
\end{abstract}

\noindent\textbf{Keywords:} Grammatical inference; Context-free languages;
Distributional learning; Positive-data identification; Finite monoids
\par\medskip

\section{Introduction}

\subsection{Background and motivation}

Gold~\cite{gold1967} showed that no language class containing all finite
languages and at least one infinite language is identifiable in the limit
from positive data alone.  For context-free languages, distributional
learning recovers positive results by restricting which substrings may be
compared across contexts; classical examples are Clark--Eyraud
substitutability~\cite{clark-eyraud2007} and Yoshinaka's
$(k,\ell)$-substitutability~\cite{yoshinaka2008}.  See
de~la~Higuera~\cite{higuera2010} for general background and
Kanazawa~\cite{kanazawa1996,kanazawa1998} for positive-data learning of
categorial grammars.  In particular, Kanazawa's monograph develops
characteristic samples for rigid grammars and finite-elasticity arguments
for the string languages of $k$-valued categorial grammars
\cite[Chapters~5 and~8]{kanazawa1998}.  This is a complementary line of
research: here the learning bias instead fixes a finite-monoid typing
of terminal factors and requires distributional substitutability.

\paragraph{Historical priority and general signature learning.}
This paper is a substantially corrected and extended version of the author's
2014 preprint~\cite{kuriyama2014}, which introduced substitutability
relative to recognizable equivalence relations and proposed positive-data
learning in that setting. The 2014 preprint predates the publication of
Yoshinaka's general many-sorted signature account~\cite{yoshinaka2015};
this chronology concerns the earlier formulation, not the validity of all
proofs or quantitative claims in that preliminary version.
In Yoshinaka's framework, for the finite-monoid homomorphism
$h:\Sigma^*\to M$ considered in this paper, the nonempty fibers
$h^{-1}(m)\cap\Sigma^+$ can be represented as disjoint sorts, with
typed concatenation and an operation promoting a typed factor to the
whole-string sort. The resulting signature has the same
substitutability condition on nonempty factors as the fixed-$h$ definition below.
Consequently, the \emph{existence} of a positive-data learner for a fixed
$h$ is also covered by the qualitative signature-based result of
\cite[Theorem~1]{yoshinaka2015}; our particular batch and sequential
constructions are analyzed directly here.
The finite-monoid-specific contributions are the algebraic fixed-window
criterion and closure-by-retyping results, the \emph{encoded size} bounds
for characteristic data (rather than only witness cardinality), and the
linear and expressive-boundary results. The broader distributional
learning setting is surveyed by Clark and Yoshinaka~\cite{clark-yoshinaka2016};
local substitutability is studied by Coste, Garet, and
Nicolas~\cite{coste-garet-nicolas2012} and Coste and Nicolas~\cite{coste-nicolas2019}.

We study a recognizable-congruence version of this program.  Let
$h:\Sigma^*\to M$ be a homomorphism into a finite monoid and
$\sim_h:=\ker h$.  For $x\in\Sigma^*$, write
$\D_L(x):=\{(u,v):uxv\in L\}$.  A language $L$ is
\emph{$\sim_h$-substitutable} if, for nonempty $x,y$,
$h(x)=h(y)$ and $\D_L(x)\cap\D_L(y)\ne\emptyset$ imply
$\D_L(x)=\D_L(y)$. Write $\Ccf{h}$ for the context-free
$\sim_h$-substitutable languages, and call this the \emph{fixed-$h$ slice}.
Thus full distributional equality is required only after a fixed finite
algebraic comparison.

For a minimal example, take $\Sigma=\{a\}$ and $L=\{a,aa\}$.
Although $(\lambda,\lambda)\in\D_L(a)\cap\D_L(aa)$,
the context $(a,\lambda)$ belongs to $\D_L(a)$ but not to $\D_L(aa)$.
Thus $L$ is not substitutable under the trivial typing.
Under the parity morphism $h(a^n)=n\bmod 2$ into
$\mathbb Z/2\mathbb Z$, however, $a$ and $aa$ have distinct types
and $L\in\Ccf{h}$.

Finite-index monoid congruences, equivalently kernels of morphisms into
finite monoids, provide finite typings.  For each fixed $h$,
write $\mathsf{RS}_h$ for the class of all $\sim_h$-substitutable
languages, without a context-freeness restriction.  For expressive
comparison, let
\[
\mathsf{RS}:=\bigcup_h\mathsf{RS}_h,
\]
where $h$ ranges over finite-monoid homomorphisms on $\Sigma^*$.
Every regular language belongs to this union.  Our learning results,
however, concern only $\Ccf{h}$ for one fixed $h$; no learner is
claimed for $\mathsf{RS}$.

Unlike one fixed prefix--suffix window, $h$ can combine independent
finite-state factor features by a product.  If the finite typing is allowed
to be refined accordingly, the framework is stable under regular filtering
and arbitrary inverse homomorphisms, including erasing ones
(Proposition~\ref{prop:finite-info-closure}).  This is a
change-of-typing statement, not closure of one fixed-$h$ slice.  An algebraic criterion identifies the typings determined by fixed
prefix--suffix windows on nonempty factors
(Proposition~\ref{prop:li-window}) and yields a corresponding
characteristic-data bound (Corollary~\ref{cor:li-thickness}).
Examples below establish nonregular linear and nonlinear reach
beyond fixed windows, and counter/Dyck obstructions bound that reach.

Specifying $h$ is a prior choice, related to typing or domain bias in
automata inference~\cite{coste-typing2004}; here types apply to entire
terminal factors, not to occurrences with their outer contexts.
For linear grammars, yield typing also induces a regular control set
of production-label sequences, giving a limited but concrete
connection with Takada's regular control-set approach
\cite{takada1988,takada1995}
(Section~\ref{sec:linear}), without an equivalence claim.


\subsection{Contributions and main results}

The paper has three main contributions.

\begin{enumerate}[label=\textup{(\arabic*)},leftmargin=2.4em]
\item \emph{Learning for fixed $h$.}  A set-driven reconstruction operator
$\mathcal B_h$ reconstructs every target in $\Ccf{h}$ exactly from finitely
many canonical positive witnesses, and a separate conservative sequential
learner $\mathcal A_h$ identifies every nonempty target in Gold's sense;
reconstruction and updates take time polynomial in the accumulated sample.

\item \emph{Quantitative guarantees.}  For arbitrary fixed $h$, we bound
the characteristic data by grammar size and by the shortest terminal yields
after target nonterminals are split according to their $h$-values.  When the
finite typing is already determined by a fixed prefix--suffix window, this
bound reduces to the ordinary shortest-yield parameter of the source grammar.
Linear fixed-$h$ targets admit characteristic data polynomial in the size of a
linear CFG.  We also show that the shortest yields after splitting can be
exponentially longer than those in the source grammar.

\item \emph{Structural and expressive position.}  We establish
closure properties obtained by finite refinement and product of typings
(rather than closure of a single fixed-$h$ slice), characterize finite
typings determined by fixed prefix--suffix windows, and relate
linear yield typing to regular control sets.  Explicit
nonregular linear and nonlinear examples lie beyond fixed windows;
counter/Dyck obstructions and comparison with Clark's congruential
family delimit the framework.
\end{enumerate}

\subsection{Organization of the paper}
Sections~\ref{sec:prelim}--\ref{sec:framework} set up the framework,
Sections~\ref{sec:algo}--\ref{sec:typing} prove learnability,
Sections~\ref{sec:typed-thickness}--\ref{sec:linear} give the quantitative
results, and Section~\ref{sec:boundaries} gives the expressiveness
comparisons.  Long normalization proofs are deferred to the appendices.

\section{Preliminaries}
\label{sec:prelim}

Let $\Sigma$ be a finite alphabet, let $\Sigma^*$ be the free monoid with
identity $\lambda$, and put $\Sigma^+:=\Sigma^*\setminus\{\lambda\}$.
A pair $(u,v)\in\Sigma^*\times\Sigma^*$ is called a
\emph{terminal context}.  For $L\subseteq\Sigma^*$ and $x\in\Sigma^*$, the \emph{distribution} of $x$ in $L$ is
$\D_L(x):=\{(u,v)\in\Sigma^*\times\Sigma^*:uxv\in L\}$.
The \emph{syntactic congruence} of $L$ is defined by $x\equiv_L y$ iff
$\D_L(x)=\D_L(y)$, for $x,y\in\Sigma^*$~\cite[Chapter~III, Proposition~10.2]{eilenberg1974}.
We say that $x$ and $y$ are \emph{syntactically congruent} when
$x\equiv_L y$.  As usual, $\equiv_L$ is a monoid congruence: if
$x\equiv_L x'$ and $y\equiv_L y'$, then $xy\equiv_L x'y'$.

For a word $w\in\Sigma^*$ and a letter $a\in\Sigma$, let $|w|_a$ denote
the number of occurrences of $a$ in $w$.

A context-free grammar (CFG) is a tuple $G=(V,\Sigma,P,S)$ with finite
nonterminal set $V$, start symbol $S\in V$, and productions
$A\to\alpha$ with $A\in V$ and $\alpha\in(V\cup\Sigma)^*$.
A \emph{sentential form} is any word in $(V\cup\Sigma)^*$.  We write
$\alpha\Rightarrow_G\beta$ for one derivation step in $G$, and
$\Rightarrow_G^*$ for its reflexive-transitive closure.  Thus $L(G):=\{w\in\Sigma^*:S\Rightarrow_G^*w\}$.  For $A\in V$, put
$L_G(A):=\{w\in\Sigma^*:A\Rightarrow_G^*w\}$.
For a nonterminal $X\in V$, a terminal context $(u,v)$ is a
\emph{terminal reaching context of $X$ in $G$} if
$S\Rightarrow_G^*uXv$.
A nonterminal is \emph{reachable} if it occurs in a sentential form reachable
from $S$, and \emph{productive} if it derives a terminal word.  A nonterminal
is \emph{useless} if it is not both reachable and productive, and it is
\emph{nullable} if it derives $\lambda$.  A production $A\to B$ with
$A,B\in V$ is a \emph{unit production} (or \emph{unit rule}).  The
\emph{unit closure} is the reflexive-transitive closure of the relation on
nonterminals induced by unit productions.  A grammar is \emph{reduced} if
every nonterminal is reachable and productive.  For a grammar generating a
nonempty language, its \emph{trim} (or \emph{reduced part}) is obtained by
first deleting all nonproductive nonterminals and their incident productions,
then deleting from the remaining grammar all nonterminals unreachable from
the start symbol, together with their incident productions.  A CFG is \emph{linear} if every right-hand side contains at
most one nonterminal.  By \emph{binarization} we mean replacement of longer
right-hand sides by equivalent chains of binary productions using fresh
nonterminals.
We write $\mathrm{CFL}$ for the class of context-free languages.  For any
finite representation $\mathcal R$, $|\mathcal R|$ denotes the length of a fixed reasonable
encoding; for a grammar this is polynomially equivalent to the usual symbol
count of its productions.  Following the thickness parameter used by
Wakatsuki--Tomita
\cite{wakatsuki-tomita1993} and adopted by Yoshinaka~\cite{yoshinaka2008}
(see also Eyraud--Heinz--Yoshinaka
\cite[Definition~2.11]{eyraud-heinz-yoshinaka2016}), for a reduced
CFG $G=(V,\Sigma,P,S)$ we use the shortest-yield \emph{thickness}
\[
 \tau_G:=\max_{A\in V}\min\{\,|w|:A\Rightarrow_G^*w,\ w\in\Sigma^*\,\}.
\]

\subsection{Learning from positive data}\label{sec:learning}

We use Gold's explanatory model~\cite{gold1967}.  Let $H_0$ denote a fixed
CFG with one start symbol and no productions, so $L(H_0)=\emptyset$.  We use
texts only for nonempty targets and represent the empty language separately by
$H_0$; equivalently one may adjoin the usual pause symbol.  For a nonempty language $L$, a \emph{text} is an
infinite sequence $T=(w_1,w_2,\ldots)$ over $L$ in which every word of $L$
occurs.  Write $T[n]=(w_1,\ldots,w_n)$ and
$K_n:=\{w_1,\ldots,w_n\}$.  A learner $\mathcal A$
\emph{identifies a nonempty $L$ in the limit} if for every text $T$ for $L$
there are $n_0$ and a grammar $H$ such that $\mathcal A(T[n])=H$ for every
$n\ge n_0$ and $L(H)=L$.
Thus the hypothesis representation itself eventually stabilizes.  A
sequential learner is \emph{conservative} if it changes its current hypothesis
only when the newly presented positive datum is not generated by that
hypothesis.

We distinguish this sequential learner from a \emph{set-driven reconstruction
operator} $\mathcal B$: this is a map from finite sets of words to grammars;
in particular, its output is independent of presentation order and multiplicity.
We call $\mathcal B$
\emph{sample-consistent} if $K\subseteq L(\mathcal B(K))$ for every finite
sample $K$.  A finite set $C\subseteq L$ is a \emph{characteristic sample of $L$ for
$\mathcal B$} if
$C\subseteq K\subseteq L$ implies $L(\mathcal B(K))=L$ for every finite
$K$.  This is the positive-data specialization of the characteristic-sample
component of de~la~Higuera's framework~\cite{higuera1997}; see also
\cite[Section~7.3.3, Definition~7.3.3]{higuera2010} for learner-relative
characteristic samples and their encoded-size convention.  The latter
is formulated for sequential presentations, whereas our condition is
explicitly set-driven.  We measure sample size by
$\|K\|:=\sum_{w\in K}(|w|+1)$,
so $\lambda$ contributes one unit.\footnote{The $+1$ prevents a nonempty sample such as $K=\{\lambda\}$ from having size $0$; it also gives $|K|\le\|K\|$, so the number of examples is controlled by the same measure.}  We use \emph{characteristic data} only
as shorthand for the encoded size $\|C\|$ of a characteristic sample $C$;
the quantitative results below combine polynomial-time reconstruction
with polynomial characteristic-data size.  Polynomial update time and
polynomial characteristic-sample size are distinct criteria in
\cite[Sections~7.3.2--7.3.3]{higuera2010}; neither alone establishes
incremental polynomial time-and-data identification.  Eyraud--Heinz--Yoshinaka
\cite[Definitions~2.5 and~2.8]{eyraud-heinz-yoshinaka2016}
distinguish incremental polynomial time-and-data identification from
its set-based counterpart.  The incremental criterion requires a correct
and stable hypothesis as soon as a characteristic sample has appeared
in the presentation; the set-based criterion requires correct
reconstruction from any finite sample containing a characteristic sample.
Our polynomial characteristic-data guarantees concern the set-driven
operator $\mathcal B_h$, not the incremental polynomial-data behavior of
$\mathcal A_h$.  Gold stabilization of the latter is established
separately in Corollary~\ref{cor:ilt}.

\subsection{Fixed finite-monoid typing}\label{sec:fixed-typing}

Fix a finite monoid $M$ and a homomorphism $h:\Sigma^*\to M$, represented
by the multiplication table of $M$ and the letter images $h(a)$ for
$a\in\Sigma$.  Thus $h(w)$ is computable by a left-to-right product in
$O(|w|)$ time; this is the sense in which the fixed homomorphism is
\emph{efficiently computable} below.

For a finite semigroup $S$, let $E(S)$ be its idempotents.  We call $S$
\emph{locally trivial} if $ete=e$ for all $e\in E(S)$ and $t\in S$.
The \emph{positive image semigroup} of $h$ is $h(\Sigma^+)$; it
need not contain the identity of $M$.

\begin{definition}[fixed-$h$ substitutability]\label{def:hsubst}
A language $L\subseteq\Sigma^*$ is \emph{$\sim_h$-substitutable} if for all
$x,y\in\Sigma^+$, equality $h(x)=h(y)$ together with
$\D_L(x)\cap\D_L(y)\ne\varnothing$ implies $\D_L(x)=\D_L(y)$.
\end{definition}
The empty word is not used as an internal substitutable fragment.  When
$\lambda\in L$, the reconstruction handles it separately through the start
symbol.

For a nonempty internal factor $x$, we call $h(x)$ its
\emph{yield type}.  As above, $\Ccf{h}$ denotes the context-free
$\sim_h$-substitutable languages, and we define
$\Clin{h}:=\{L\in\Ccf{h}:L\text{ is generated by some linear CFG}\}$.

\section{Fixed Recognizable Substitutability}
\label{sec:framework}

For fixed $h$, put
$\mathsf{RS}_h:=\{L\subseteq\Sigma^*:L\text{ is $\sim_h$-substitutable}\}$.
Then the fixed-$h$ slice is
$\Ccf{h}=\mathsf{RS}_h\cap\mathrm{CFL}$, while $\Clin{h}$ consists
of its linearly generated members.  For expressive comparisons only, let
$\mathsf{RS}:=\bigcup_h\mathsf{RS}_h$, where $h$ ranges over finite-monoid
homomorphisms on $\Sigma^*$.  We write $\mathsf{RS}_h(\Gamma)$ and
$\mathsf{RS}(\Gamma)$ when the alphabet matters; all inclusions below are
alphabetwise.  The learning results always concern one fixed $h$.

\begin{proposition}\label{prop:regular-auto}
A language is regular iff it has the form $h^{-1}(\mathsf{Acc})$ for some
homomorphism $h:\Sigma^*\to M$ into a finite monoid and some
$\mathsf{Acc}\subseteq M$.  For regular $L$, $h$ may be taken to be its
syntactic morphism, whose kernel is $\equiv_L$.  For every fixed $h$ and
$\mathsf{Acc}\subseteq M$, the language $h^{-1}(\mathsf{Acc})$ is
$\sim_h$-substitutable.  Thus every regular language belongs to
$\mathsf{RS}$.
\end{proposition}

\begin{proof}
The first two statements follow from the classical syntactic-monoid
characterization~\cite[Chapter~III, Propositions~10.1--10.2]{eilenberg1974}.  If $L=h^{-1}(\mathsf{Acc})$ and $h(x)=h(y)$, then
$h(uxv)=h(uyv)$ for all $u,v$, so $\D_L(x)=\D_L(y)$.
\end{proof}

By Proposition~\ref{prop:regular-auto}, $\mathsf{RS}$ and
$\mathsf{RS}\cap\mathrm{CFL}$ contain every finite language and
$\Sigma^*$; Gold's theorem makes both unions nonlearnable from
positive data on nonempty alphabets.  Yoshinaka
\cite[Section~4]{yoshinaka2008} notes the analogous issue for
the union of fixed-window classes.  We learn only fixed-$h$ slices.

Typing refinement is monotone on the nonempty factors that enter
Definition~\ref{def:hsubst}.  If
$\ker h'|_{\Sigma^+}\subseteq\ker h|_{\Sigma^+}$, then
$\mathsf{RS}_h\subseteq\mathsf{RS}_{h'}$, since $h'$ requires fewer
factor comparisons.  Refinement may nevertheless enlarge the typed grammar
used in the quantitative analysis.  For two finite typings $h,g$,
$\ker(h\times g)|_{\Sigma^+}
 =\ker h|_{\Sigma^+}\cap\ker g|_{\Sigma^+}$, hence
$\mathsf{RS}_h\cup\mathsf{RS}_g\subseteq\mathsf{RS}_{h\times g}$;
in particular, window information and another finite feature can be combined
without making the typing target-specific.

Yoshinaka~\cite[Proposition~1]{yoshinaka2008} showed that
fixed-window substitutable languages are not closed under intersection
with regular languages or inverse images under homomorphisms that may erase.
In contrast, they are closed under inverse images of nonerasing
homomorphisms.  The following proposition shows that regular filtering
and erasing inverse images can be accommodated by changing the finite typing.

\begin{proposition}[closure under finite information]\label{prop:finite-info-closure}
Let $h:\Sigma^*\to M$ and $g:\Sigma^*\to N$ be finite-monoid
homomorphisms.
\begin{enumerate}[label=\textup{(\roman*)},nosep]
\item If $L\in\mathsf{RS}_h$ and $J\in\mathsf{RS}_g$, then
$L\cap J\in\mathsf{RS}_{h\times g}$.
\item If $Q=g^{-1}(F)$ is regular and $L\in\Ccf{h}$, then
$L\cap Q\in\Ccf{h\times g}$.
\item Let $\varphi:\Gamma^*\to\Sigma^*$ be any homomorphism, possibly
erasing.  Let $B_\varphi=\{1,\zeta\}$ be the two-element monoid with zero
$\zeta$, define $e_\varphi(w)=1$ iff $\varphi(w)=\lambda$ and
$e_\varphi(w)=\zeta$ otherwise, and put
$\widehat h=(h\circ\varphi)\times e_\varphi$.
If $L\in\mathsf{RS}_h$, then
$\varphi^{-1}(L)\in\mathsf{RS}_{\widehat h}$; if moreover
$L\in\Ccf{h}$, then
$\varphi^{-1}(L)\in\Ccf{\widehat h}$.
\end{enumerate}
\end{proposition}

\begin{proof}
For (i), suppose nonempty $x,y$ have equal $(h\times g)$-type and share a
context in $L\cap J$.  They then have equal $h$-type and a common context in
$L$, and equal $g$-type and a common context in $J$.  Hence
$\D_L(x)=\D_L(y)$ and $\D_J(x)=\D_J(y)$, so
$\D_{L\cap J}(x)=\D_{L\cap J}(y)$.  Item~(ii) follows from (i),
Proposition~\ref{prop:regular-auto}, and closure of CFLs under intersection
with regular languages.

For (iii), $e_\varphi$ is a homomorphism because
$\varphi(uv)=\lambda$ iff both $\varphi(u)$ and $\varphi(v)$ are empty.
Let equal-$\widehat h$-type nonempty $x,y$ share a context in
$\varphi^{-1}(L)$.  If $e_\varphi(x)=e_\varphi(y)=1$, then
$\varphi(x)=\varphi(y)=\lambda$, so the two factors have identical
distributions in the inverse image.  Otherwise both images are nonempty,
have equal $h$-type, and the shared inverse-image context maps to a common
context of $\varphi(x),\varphi(y)$ in $L$.  Fixed-$h$ substitutability in
$L$ gives equal distributions there, and pulling contexts back through
$\varphi$ gives
$\D_{\varphi^{-1}(L)}(x)=\D_{\varphi^{-1}(L)}(y)$.
The context-free statement uses closure of CFLs under inverse homomorphism.
\end{proof}

\paragraph{When is a typing useful?}
For correctness relative to a target $L$, Definition~\ref{def:hsubst} says
exactly that $h$ must separate every pair of nonempty factors that share a
context but have different distributions.  Equivalently, choosing a useful
typing is a finite-index monoid-congruence separation problem: no such
distributional conflict may be collapsed by $\ker h$.  Typing refinement is
therefore safe, and product typing combines independent finite features because
a conflicting pair need only be separated by one component.  Lemma~\ref{lem:finite-monoid-obstruction}
later gives the corresponding obstruction: an infinite pairwise family of
conflicting factors cannot be separated by any finite typing.  Quantitative
usefulness is a separate issue.  Refinement can enlarge the typed grammar used
by the reconstruction bound; Proposition~\ref{prop:typed-thickness-gap} exhibits an
exponential source-to-typed thickness gap, whereas
Sections~\ref{sec:fixed-window-thick} and~\ref{sec:linear} give sufficient
regimes in which that cost is controlled.

If $\ker h$ restricted to $\Sigma^+$ refines the syntactic congruence
$\equiv_L$, then $L\in\mathsf{RS}_h$.  The converse already fails for
Clark--Eyraud substitutability.  Consider
$L_{\mathrm c}:=\{a^ncb^n:n\ge0\}$.  Under the nonempty-fragment
convention of Definition~\ref{def:hsubst}, $L_{\mathrm c}$ is Clark--Eyraud substitutable,
but $\equiv_{L_{\mathrm c}}$ has infinite index: for $1\le i<j$ the
context $(a,cb^{i+1})$ distinguishes $a^i$ from $a^j$.


\subsection{Classical special cases}

Clark--Eyraud's original Definition~5 includes all fragments
\cite{clark-eyraud2007}; we follow Clark~\cite{clark2013}
in excluding $\lambda$.  This distinction matters even for $L=a^+$:
$(a,\lambda)$ belongs to both $\D_L(\lambda)$ and $\D_L(a)$,
whereas $(\lambda,\lambda)$ belongs to $\D_L(a)$ but not to
$\D_L(\lambda)$.  Thus comparing $\lambda$ with nonempty factors
would exclude $a^+$.
Call $L$
\emph{Clark--Eyraud substitutable} when
$\D_L(x)\cap\D_L(y)\ne\emptyset$ implies
$\D_L(x)=\D_L(y)$ for all $x,y\in\Sigma^+$.  This is exactly the
trivial-monoid case of Definition~\ref{def:hsubst}, and hence also
Yoshinaka's $(0,0)$ case.


For $k,\ell\ge0$, let $\pref_k(w)$ and $\suf_\ell(w)$ denote the
prefix and suffix of lengths $\min\{k,|w|\}$ and
$\min\{\ell,|w|\}$, respectively; in particular,
$\pref_0(w)=\suf_0(w)=\lambda$.  Renaming Yoshinaka's variables so that
$u_i,v_i$ denote outer contexts, $L$ is \emph{$(k,\ell)$-substitutable}
if, for all $p\in\Sigma^k$, $q\in\Sigma^\ell$, and
$u_1,u_2,v_1,v_2,r,s\in\Sigma^*$ with $prq,psq\ne\lambda$,
\[
\begin{split}
&u_1prqv_1\in L,\quad u_1psqv_1\in L,\quad
u_2prqv_2\in L\\
&\hspace{35mm}\Longrightarrow\quad u_2psqv_2\in L.
\end{split}
\]

\begin{proposition}\label{prop:yl-special}
For every fixed $k,\ell\ge0$ there is a finite monoid homomorphism
$h_{k,\ell}$ such that, for every $L\subseteq\Sigma^*$,
$L$ is $(k,\ell)$-substitutable iff it is
$\sim_{h_{k,\ell}}$-substitutable.
\end{proposition}

\begin{proof}
Put $m_0:=\max\{1,k+\ell\}$, let $h_{k,\ell}(w):=w$ for
$|w|<m_0$, and otherwise let
$h_{k,\ell}(w):=\langle\pref_k(w),\suf_\ell(w)\rangle$, with a tag
separating the two cases.  Denote the finite image by
$M_{k,\ell}:=h_{k,\ell}(\Sigma^*)$.  On $M_{k,\ell}$ define, for
$x,y$ with $|x|,|y|<m_0$,
\[
\begin{aligned}
x\cdot y&:=h_{k,\ell}(xy),&
x\cdot\langle p,q\rangle&:=\langle\pref_k(xp),q\rangle,\\
\langle p,q\rangle\cdot y&:=\langle p,\suf_\ell(qy)\rangle,&
\langle p,q\rangle\cdot\langle p',q'\rangle&:=\langle p,q'\rangle.
\end{aligned}
\]
Then $h_{k,\ell}(xy)=h_{k,\ell}(x)\cdot h_{k,\ell}(y)$; surjectivity
gives associativity and $h_{k,\ell}(\lambda)$ is the identity.

For nonempty $x,y$, equal types mean either $x=y$ or
$x=prq$, $y=psq$ for a common length-$k$ prefix $p$ and length-$\ell$
suffix $q$ (including $k=\ell=0$).  If $L$ is
$(k,\ell)$-substitutable, a shared context plus any context of one factor
instantiates Yoshinaka's four-word condition, and symmetry gives equal
distributions.  Conversely, in any Yoshinaka premise the factors
$x:=prq$ and $y:=psq$ are nonempty and lie in the long-word part of the
coding, hence have equal $h_{k,\ell}$-type; fixed-$h_{k,\ell}$
substitutability transfers the third context and yields the fourth word.
\end{proof}

\phantomsection\label{crit:kl}
\paragraph{Fixed-window counterexample criterion.}
If $x,y$ have length at least $\max\{1,k+\ell\}$, the same length-$k$
prefix and length-$\ell$ suffix, a common context, and unequal
distributions, then $L$ is not $(k,\ell)$-substitutable: the first
three conditions give equal $h_{k,\ell}$-types and a shared context, while
the last contradicts $\sim_{h_{k,\ell}}$-substitutability.

Thus the fixed $(k,\ell)$ class is exactly
$\mathsf{RS}_{h_{k,\ell}}$.  Section~\ref{sec:fixed-window-thick} gives
its polynomial thick-data bound.  For $k=\ell=0$,
$h_{0,0}$ merely separates $\lambda$ from nonempty words and induces the
same substitutability condition as the one-element monoid.  For an alphabet
$\Gamma$, set
\[
\mathsf{KL}(\Gamma):=
\bigcup_{k,\ell\ge0}
\{L\subseteq\Gamma^*:L\text{ is $(k,\ell)$-substitutable}\},
\]
and abbreviate $\mathsf{KL}:=\mathsf{KL}(\Sigma)$.

\subsection{Main theorem}

\begin{theorem}[fixed-$h$ learning theorem]\label{thm:main}
Fix an efficiently computable homomorphism $h:\Sigma^*\to M$ into a finite
monoid.  There are a set-driven reconstruction operator $\mathcal B_h$ and a
conservative sequential learner $\mathcal A_h$ such that:
\begin{enumerate}[label=\textup{(\roman*)},nosep]
\item $\mathcal B_h$ is sample-consistent and $\mathcal B_h(K)$ is
constructible in time polynomial in $\|K\|$;
\item every $L\in\Ccf{h}$ has a finite characteristic sample for
$\mathcal B_h$;
\item $\mathcal A_h$ identifies every nonempty $L\in\Ccf{h}$ in Gold's
sense, each update taking time polynomial in the data seen so far;
\item if the positive image semigroup $h(\Sigma^+)$ is locally trivial,
then for every nonempty target the sample in~(ii) can be chosen polynomial
in $|G_*|$ and $\tau_{G_*}$ for any reduced CFG $G_*$ representing that
target; in particular this applies to every
fixed-window typing $h_{k,\ell}$ (also directly by
Theorem~\ref{thm:window-thick} and
Proposition~\ref{prop:thick-ssbnf-normal});
\item for arbitrary fixed $h$ and $L\in\Clin{h}$, the sample in~(ii) can
be chosen polynomial in $|G_*|$ for any linear CFG $G_*$ representing $L$.
\end{enumerate}
The polynomials in~(iv)--(v) may depend on the fixed $h$ and $\Sigma$,
but not on the chosen target grammar $G_*$. The quantitative mechanism behind items~(iv)--(v), together with its
limitation for arbitrary fixed typings, is developed in
Sections~\ref{sec:typed-thickness}--\ref{sec:linear}.
\end{theorem}

This is consistent with Gold's theorem: on a nonempty alphabet choose
distinct $x,y\in\Sigma^+$ with $h(x)=h(y)$, and choose
$z\in\Sigma^*$ with $|z|>\max\{|x|,|y|\}$.  The finite language $\{x,y,zx\}$
is not $\sim_h$-substitutable: $x,y$ share the empty context
but only $x$ admits $(z,\lambda)$.

\paragraph{Algebraic position of fixed windows.}
The following positive-image/kernel criterion is classical
\cite[Propositions~XI.4.17 and~XIV.1.20]{pin2025};
the generalized definite language characterization is likewise known
\cite[Proposition~5.18]{pin1997}.
Its learning-specific application is Corollary~\ref{cor:li-thickness}.

\begin{proposition}[locally trivial typings and fixed windows]
\label{prop:li-window}
Let $h:\Sigma^*\to M$ be a finite-monoid homomorphism on a nonempty
alphabet, and put $n:=|h(\Sigma^+)|+1$.  The following are equivalent:
\begin{enumerate}[label=\textup{(\roman*)},nosep]
\item $h(\Sigma^+)$ is locally trivial;
\item there are $k,\ell\ge0$ such that
$\ker h_{k,\ell}|_{\Sigma^+}\subseteq\ker h|_{\Sigma^+}$.
\end{enumerate}
If these conditions hold, $(k,\ell)=(n,n)$ works.  Moreover every positive
image semigroup $h_{k,\ell}(\Sigma^+)$ is locally trivial, and therefore
\[
 \mathsf{KL}(\Sigma)
 =
 \bigcup_{\substack{h:\Sigma^*\to M\\ h(\Sigma^+)\text{ locally trivial}}}
 \mathsf{RS}_h(\Sigma).
\]
The same equality holds after intersecting both sides with $\mathrm{CFL}$.
\end{proposition}

\begin{proof}
Assume (i) and put $S=h(\Sigma^+)$, $n=|S|+1$. Then
$S^n=SE(S)S$, and local triviality also implies $esf=ef$ for
all $e,f\in E(S)$ and every $s\in S$
\cite[Corollary~II.6.35 and Proposition~XI.4.17]{pin2025}.
Thus for $p,q\in\Sigma^n$
we may write $h(p)=aeb$, $h(q)=cfd$, with $e,f$ idempotent, and obtain
$h(prq)=a\,e\,(b\,h(r)\,c)\,f\,d=aefd=h(pq)$
for every $r\in\Sigma^*$. The definition of $h_{n,n}$, which
distinguishes short positive words individually, gives (ii).

Conversely, if $S$ is not locally trivial, choose an idempotent
$e\in S$ and $m\in S$ with $eme\ne e$, represented by nonempty
words $r,z$, respectively. Given $k,\ell$, take
$j\ge1$ with $j|r|\ge\max\{k,\ell\}$.
Then $x=r^{2j+1}$ and $y=r^jzr^j$ have the same
$h_{k,\ell}$-value, but $h(x)=e\ne eme=h(y)$. Hence (ii) fails.

Finally, the multiplication in Proposition~\ref{prop:yl-special}
shows that every idempotent of the positive image of $h_{k,\ell}$
has the form $\langle p,q\rangle$ and satisfies
$\langle p,q\rangle t\langle p,q\rangle=\langle p,q\rangle$
for every positive-image element $t$; thus that semigroup is locally
trivial. The forward implication and the kernel-refinement monotonicity
give $\mathsf{RS}_h\subseteq\mathsf{RS}_{h_{n,n}}$ whenever
$h(\Sigma^+)$ is locally trivial. Taking unions yields the stated
equality, also after restriction to $\mathrm{CFL}$.
\end{proof}
\section{Reconstruction Operator and Sequential Learner}
\label{sec:algo}


\subsection{Finite-sample reconstruction}
For a finite sample $K\subseteq\Sigma^*$, let
\[
 \widehat V(K):=
 \{[x]:x\in\Sigma^+,\ \exists u,v\in\Sigma^*\text{ with }uxv\in K\}.
\]
Thus there is one hypothesis nonterminal for each distinct nonempty factor
observed in the sample; the brackets are part of the nonterminal name and do
not denote a quotient class.  For observed nonempty factors $x,y$, write
\[
 x\approx_{K,h}y
 \quad\Longleftrightarrow\quad
 h(x)=h(y)\ \text{ and }\ \D_K(x)\cap\D_K(y)\ne\emptyset .
\]
The grammar $\mathcal B_h(K)$ has nonterminal set
$\widehat V(K)\cup\{\widehat S\}$, start symbol $\widehat S$, and the following
productions whenever every displayed bracketed nonterminal belongs to
$\widehat V(K)$:
\[
\begin{array}{rll}
\textup{(B)}&[xy]\to[x][y],&x,y\in\Sigma^+,\\[1mm]
\textup{(U)}&[x]\to[y],&x\approx_{K,h}y,\\[1mm]
\textup{(L)}&[a]\to a,&a\in\Sigma,\\[1mm]
\textup{(S)}&\widehat S\to[w],&w\in K\cap\Sigma^+.
\end{array}
\]
If $\lambda\in K$, Rule~\textup{(S)} also contains
$\widehat S\to\lambda$.

The construction follows the substring-indexed grammar scheme of
Yoshinaka~\cite[Section~4, Algorithm~1]{yoshinaka2008}, also recalled
by Clark~\cite[Section~4, Algorithm~1]{clark2013}.
For $h=h_{k,\ell}$, the nontrivial instances of Rule~\textup{(U)}
coincide with Yoshinaka's unary rules: distinct observed nonempty factors
are connected when they share a sample context and have equal
$h_{k,\ell}$-type; equivalently, both factors have length at least
$\max\{1,k+\ell\}$ and have the same length-$k$ prefix and
length-$\ell$ suffix.
The constructions agree up to redundant identity rules, the
representation of initial productions by a separate start symbol, and
the treatment of $\lambda$.  For general fixed $h$, the
prefix--suffix test is replaced by equality of $h$-types.

\begin{lemma}[sample consistency]\label{lem:sample-consistency}
For every finite sample $K$, $K\subseteq L(\mathcal B_h(K))$.
\end{lemma}
\begin{proof}
If $w=a_1\cdots a_m\in K\cap\Sigma^+$, Rule~\textup{(S)} enters $[w]$.
Repeated use of Rule~\textup{(B)} along any binary factorization of $w$
reaches the one-letter factors, all of which are observed because they occur
inside the same sample word.  Rule~\textup{(L)} then emits
$a_1\cdots a_m$.  If $\lambda\in K$, the start rule emits it directly.
\end{proof}

\begin{theorem}[soundness]\label{thm:soundness}
If $K\subseteq L$ and $L$ is $\sim_h$-substitutable, then
$L(\mathcal B_h(K))\subseteq L$.
\end{theorem}

\begin{proof}
Induct on the height of a derivation tree rooted at $[x]$ to prove
\begin{equation}\label{eq:soundness-invariant}
 [x]\Rightarrow_{\mathcal B_h(K)}^*w
 \quad\Longrightarrow\quad
 x\equiv_L w\ \text{ and }\ h(w)=h(x).
\end{equation}
Every factor nonterminal is indexed by a nonempty observed word,
so the fragments compared by Rule~\textup{(U)} satisfy the nonempty
requirement of Definition~\ref{def:hsubst}.  Moreover,
Rules~\textup{(B)}--\textup{(L)} cannot derive $\lambda$ from a
factor nonterminal.  Rule~\textup{(L)} is immediate.  For Rule~\textup{(U)}, suppose
$[x]\to[y]$.  By definition there is a sample context $(u,v)$ with
$uxv,uyv\in K\subseteq L$ and $h(x)=h(y)$.  Fixed-$h$ substitutability gives
$x\equiv_L y$; the induction hypothesis gives $y\equiv_L w$ and
$h(w)=h(y)$, proving~\eqref{eq:soundness-invariant}.

For Rule~\textup{(B)}, write the derived word as $w_1w_2$ from
$[xy]\to[x][y]$.  The induction hypotheses give
$x\equiv_L w_1$, $y\equiv_L w_2$,
$h(w_1)=h(x)$, and $h(w_2)=h(y)$.
Because $\equiv_L$ is a monoid congruence,
$xy\equiv_L w_1w_2$, and the homomorphism law gives
$h(w_1w_2)=h(xy)$.

Finally, if a start rule $\widehat S\to[x]$ is used, then $x\in K\subseteq L$,
so $(\lambda,\lambda)\in\D_L(x)$.  The invariant gives
$\D_L(w)=\D_L(x)$ and hence $w\in L$.  The direct $\lambda$ rule is present
only when $\lambda\in K\subseteq L$.
\end{proof}

\subsection{The sequential learner}

Recomputing $\mathcal B_h$ after every datum need not stabilize
syntactically: for $L_{\mathrm c}=\{a^ncb^n:n\ge0\}$, longer examples
introduce new factor nonterminals such as $[a^ncb^n]$.
We therefore use the conservative wrapper of Clark--Eyraud's
Substitution Graph Learner (SGL, Algorithm~2) and Yoshinaka's
$(k,\ell)$-Substitutable Grammar Learner ($(k,\ell)$-SGL, Algorithm~1)~\cite{clark-eyraud2007,yoshinaka2008};
see also Eyraud--Heinz--Yoshinaka~\cite[Appendix]{eyraud-heinz-yoshinaka2016}
for the distinction between set-driven reconstruction and incremental learning.
Starting from the fixed empty grammar $H_0$, after reading $w_n$ put
\[
 H_n:=
 \begin{cases}
 H_{n-1},&w_n\in L(H_{n-1}),\\
 \mathcal B_h(K_n),&w_n\notin L(H_{n-1}),
 \end{cases}
 \qquad
 \mathcal A_h(T[n]):=H_n.
\]
Every rebuild uses the whole accumulated sample $K_n$.  Although
$\mathcal B_h$ is set-driven, the grammar on which $\mathcal A_h$ stabilizes
may depend on the presentation order; only eventual syntactic stabilization
to a grammar for the target is claimed (Corollary~\ref{cor:ilt}).
This order dependence is genuine: for any infinite target and a
convergent text $T$, choose $w\in L$ longer than all examples at
the last rebuild.  The limiting grammar for $T$ omits $[w]$,
but prepending $w$ to $T$ puts $[w]$ in every rebuilt grammar,
so the two limiting grammars differ.
Clark~\cite[Section~4, Example~1]{clark2013} illustrates the same
presentation-order dependence for
$L_{\mathrm c}=\{a^ncb^n:n\ge0\}$: different initial positive
examples can lead a conservative learner to stabilize on distinct
substring-indexed grammars.

\section{Finite Witnesses, Exact Reconstruction, and Complexity}
\label{sec:typing}


\subsection{Target refinement and finite witnesses}

For the completeness proof only, we refine a target grammar by yield type and
show that the hypothesis can mimic finitely witnessed typed derivations.  We use
the following normal form as a proof device.

\begin{definition}[start-separated binary normal form]
A CFG $G=(V,\Sigma,P,S_0)$ is in \emph{start-separated binary normal form}
(SSBNF) if $S_0$ never occurs on a right-hand side,
every start rule is $S_0\to A$ or $S_0\to\lambda$, and every non-start rule
is either $A\to a$ or $A\to BC$.
\end{definition}
The start separation isolates whole-sentence entry from internal
nonterminals and allows $\lambda$ to be handled at the start symbol.

Every nonempty context-free language has an equivalent reduced SSBNF grammar by
the standard start-separation, useless-symbol elimination, $\lambda$/unit elimination
away from the start, and binarization procedures.  The refinement below is
therefore defined for nonempty context-free languages; the empty target is
represented separately by $H_0$ and $\mathcal B_h(\emptyset)$.

Let $G=(V,\Sigma,P,S_0)$ be a reduced SSBNF grammar for a nonempty language
$L$.  Form the following yield-typed grammar with start symbol
$\widetilde S$ and non-start symbols $A_\mu$ with
$A\in V\setminus\{S_0\}$ and $\mu\in M$.  Here and below, juxtaposition of type indices denotes multiplication in
$M$; thus $\mu\nu=\mu\cdot\nu$.  The rules are
\[
\begin{array}{ll}
\widetilde S\to A_\mu & (S_0\to A\in P,\ \mu\in M),\\
\widetilde S\to\lambda & (S_0\to\lambda\in P),\\
A_{h(a)}\to a & (A\to a\in P),\\
A_{\mu\nu}\to B_\mu C_\nu & (A\to BC\in P,\ \mu,\nu\in M).
\end{array}
\]
The reduced part of this grammar is denoted
$\widetilde G=(\widetilde V,\Sigma,\widetilde P,\widetilde S)$ and called
the \emph{typed refinement}.  By construction, every symbol of
$\widetilde G$ is therefore reachable and productive.

Fix an order on $\Sigma$ and use the induced \emph{shortlex} order:
shorter strings come first, with lexicographic order breaking equal-length
ties.  For every non-start symbol $X$ of $\widetilde G$, let $\omega(X)$
be its shortlex-minimal terminal yield, called the \emph{canonical yield} of
$X$.  This is well-defined because $\widetilde G$ is reduced.

The yield-type split is used in the present completeness proof: the canonical
yield of a typed parent and the concatenated canonical yields of
its typed children must have equal $h$-types to license unary
Rule~\textup{(U)} (Theorem~\ref{thm:complete}).

The role of this split in the present completeness argument can already be seen
before typing.  Let
$h(w)=|w|\bmod 2$ and let $G$ be the SSBNF grammar with start symbol $S_0$
and rules
$S_0\to A$, $A\to a\mid BC$, $B\to a$, and $C\to a$.  The shortest yield of $A$
is $a$, of odd type, while the concatenation of shortest yields of
$B$ and $C$ is $aa$, of even type.  Without separating $A$ by yield type,
the corresponding unary-rule step is therefore not licensed.
Splitting by yield type makes the required equality an invariant.  By
contrast, no types are attached to the surrounding left or right contexts,
because the reconstruction operator never tests such types.

\begin{proposition}[typed lifting and yield invariant]\label{prop:typed-core}
The typed refinement satisfies $L(\widetilde G)=L(G)$.  For every retained
non-start symbol $A_\mu$,
\[
L_{\widetilde G}(A_\mu)=L_G(A)\cap h^{-1}(\{\mu\}).
\]
In particular, $A_\mu\Rightarrow_{\widetilde G}^*w$ implies $h(w)=\mu$.
\end{proposition}
\begin{proof}
Erase type annotations to map every typed derivation to a derivation of $G$.
Conversely, annotate each non-start node of a successful derivation tree of
$G$ by the $h$-type of its terminal yield.  The terminal and binary rules
above are exactly the bottom-up propagation laws for those types, so the
annotated tree is a derivation of the untrimmed yield-typed grammar above;
all symbols on a successful tree survive trimming.

For a retained non-start $A_\mu$, erasing annotations proves
$L_{\widetilde G}(A_\mu)\subseteq L_G(A)\cap h^{-1}(\{\mu\})$.
Conversely, fix $w\in L_G(A)$ with $h(w)=\mu$.
Since $A_\mu$ survives trimming, there is a successful typed derivation
tree containing an occurrence of $A_\mu$.  Annotate a source derivation
$A\Rightarrow_G^*w$ bottom-up by the $h$-types of its subtree yields and
graft the resulting typed subtree at that occurrence.  Its root has type
$\mu$, so the graft is valid; the resulting global typed tree is successful.
Consequently all symbols and rules of the grafted subtree survive trimming,
and $A_\mu\Rightarrow_{\widetilde G}^*w$.
The yield-type invariant also follows directly by induction on a typed
derivation.
\end{proof}

For every non-start symbol $X$ of $\widetilde G$, choose
$\chi(X)=(u_X,v_X)$ among the terminal reaching contexts of $X$ so as to
minimize $|u_X|+|v_X|$, breaking ties by the shortlex orders of $u_X$ and
then $v_X$; call $\chi(X)$ the \emph{canonical context} of $X$.  Such a
terminal reaching context exists because $\widetilde G$ is reduced:
reachability places $X$ in a reachable sentential form, and productivity of
the surrounding nonterminals lets all symbols to its left and right be
expanded to terminal strings.  Proposition~\ref{prop:typed-core} gives
$h(\omega(A_\mu))=\mu$ for every typed non-start symbol.

Define the finite witness set
\begin{equation}\label{eq:W}
\begin{split}
\WitnessSet(\widetilde G):={}&
 \{u_X\omega(X)v_X:X\in\widetilde V\setminus\{\widetilde S\}\}\\
&\cup\{u_Xav_X:X\to a\in\widetilde P\}\\
&\cup\{u_X\omega(Y)\omega(Z)v_X:X\to YZ\in\widetilde P\}\\
&\cup\bigl(\{\lambda\}\text{ if }\lambda\in L\bigr).
\end{split}
\end{equation}
Every displayed word belongs to $L$, and the set is finite.  In particular,
for a binary rule $X\to YZ$ the word in the third family is justified by
$\widetilde S\Rightarrow^*u_XXv_X\Rightarrow u_XYZv_X
\Rightarrow^*u_X\omega(Y)\omega(Z)v_X$; the first and second families are
justified analogously.  We call the elements of $\WitnessSet(\widetilde G)$
the \emph{canonical positive witnesses}.  The first family records one
canonical occurrence of each typed symbol, and the second and third families
record the terminal and binary productions, respectively.


\subsection{Exact reconstruction and Gold convergence}

\begin{theorem}[completeness from the finite witnesses]\label{thm:complete}
Let $G$ be any reduced SSBNF grammar for a nonempty language $L$, and let
$\widetilde G$ be its typed refinement with respect to the fixed
homomorphism $h$.  If a finite set $K$ satisfies
$\WitnessSet(\widetilde G)\subseteq K$, then $L\subseteq L(\mathcal B_h(K))$.
\end{theorem}

Neither substitutability nor $K\subseteq L$ is needed for
this combinatorial completeness statement.
\begin{proof}
For each non-start $X$ write $\widehat X:=[\omega(X)]$.
The corresponding word in the first family of~\eqref{eq:W} ensures that
$\widehat X$ exists.  We prove by induction on a typed derivation
$X\Rightarrow^*w$ that
$\widehat X\Rightarrow_{\mathcal B_h(K)}^*w$.

If the first rule is $X\to a$, the first and second witness families contain
$u_X\omega(X)v_X$ and $u_Xav_X$, respectively.  Hence
$\omega(X)\approx_{K,h}a$ because
$h(\omega(X))=h(a)$ by Proposition~\ref{prop:typed-core}.  Rule~\textup{(U)}
moves $\widehat X$ to $[a]$ when necessary, and Rule~\textup{(L)} yields $a$.

Suppose the first rule is $X=A_{\mu\nu}\to B_\mu C_\nu$; put
$Y:=B_\mu$ and $Z:=C_\nu$.  The first and third witness families contain
$u_X\omega(X)v_X$ and
$u_X\omega(Y)\omega(Z)v_X$.  Since
\[
 h(\omega(X))=\mu\nu=h(\omega(Y))h(\omega(Z))
              =h(\omega(Y)\omega(Z)),
\]
Rule~\textup{(U)} moves $\widehat X$ to
$[\omega(Y)\omega(Z)]$ when necessary.  The third witness also makes
$\omega(Y)\omega(Z)$ an observed factor, so Rule~\textup{(B)} gives
\[
 [\omega(Y)\omega(Z)]\to[\omega(Y)][\omega(Z)].
\]
The induction hypotheses finish the two subderivations.

Finally, every nonempty $w\in L$ has a lifted start derivation
$\widetilde S\to X\Rightarrow^*w$.  The start rule itself shows that
$(\lambda,\lambda)$ is a terminal reaching context of $X$; its total length is
$0$, which is minimal, so the canonical choice is
$\chi(X)=(\lambda,\lambda)$.  The first witness family therefore gives
$\omega(X)\in K$, and Rule~\textup{(S)} starts the preceding derivation.
The case $w=\lambda$ is immediate from~\eqref{eq:W} and Rule~\textup{(S)}.
\end{proof}

\begin{theorem}[exact finite-sample reconstruction]
\label{thm:reconstruction-fixed-h}
Let $L\in\Ccf{h}$ be nonempty, let $G$ be a reduced SSBNF grammar for $L$,
and let $\widetilde G$ be its typed refinement.  If
$\WitnessSet(\widetilde G)\subseteq K\subseteq L$, then $L(\mathcal B_h(K))=L$.
\end{theorem}
\begin{proof}
Combine Theorems~\ref{thm:soundness} and~\ref{thm:complete}.
\end{proof}

Thus $\WitnessSet(\widetilde G)$ is characteristic for every
nonempty target, including $\{\lambda\}$.  For $L=\emptyset$,
the empty sample gives $L(\mathcal B_h(\emptyset))=\emptyset$.
This proves Theorem~\ref{thm:main}(ii).

\begin{corollary}[Gold identification]\label{cor:ilt}
For a fixed homomorphism $h:\Sigma^*\to M$ into a finite monoid $M$, the learner
$\mathcal A_h$ identifies every nonempty language in $\Ccf{h}$ from positive data.  More precisely, for any reduced SSBNF target
grammar $G$ and its typed refinement $\widetilde G$, once
$\WitnessSet(\widetilde G)\subseteq K_{n_0}$, among stages $n\ge n_0$ at most one hypothesis change
can occur.
\end{corollary}

\begin{proof}
Under our text convention the Gold statement concerns nonempty targets;
$H_0$ is the initial empty hypothesis.  Let $L$ be nonempty and let $T$ be a
text for $L$.  Since $w_1\notin L(H_0)=\emptyset$, the learner rebuilds at
stage $1$.  Choose $n_0$ with
$\WitnessSet(\widetilde G)\subseteq K_{n_0}$.  If the learner rebuilds at
some stage $n\ge n_0$, Theorem~\ref{thm:reconstruction-fixed-h} makes the
new hypothesis exactly $L$, after which no positive example can trigger
another change.  In this case that rebuild is the last hypothesis change.
Suppose instead that no rebuild occurs at any stage $n\ge n_0$.  Soundness
gives $L(H_{n_0})\subseteq L$.  Moreover, $K_n\subseteq L(H_n)$ for every
$n$: a rebuild is sample-consistent, while at a non-rebuild step the newly
presented datum is already generated by the unchanged hypothesis.  Hence a
strict inclusion $L(H_{n_0})\subsetneq L$ would leave some
$w^\star\in L\setminus L(H_{n_0})$ that has not yet occurred by stage
$n_0$.  Since every target word eventually occurs in the text, $w^\star$
must occur later and force a rebuild, a contradiction.  Thus
$L(H_{n_0})=L$; in this case no hypothesis change after stage $n_0$ is
needed.
\end{proof}

A missed-positive rebuild can be delayed after a characteristic
sample appears; compare Eyraud--Heinz--Yoshinaka
\cite[Appendix]{eyraud-heinz-yoshinaka2016}.

\subsection{Complexity of the reconstruction step}
\label{sec:complexity}

\begin{theorem}[polynomial-time reconstruction]\label{thm:poly-build}
For fixed $h$, $\mathcal B_h(K)$ is constructible in time polynomial in
$\|K\|$.
\end{theorem}
\begin{proof}
Put $n_K:=\|K\|$.  Across all sample words there are $O(n_K^2)$ nonempty
factor occurrences, hence at most $O(n_K^2)$ distinct factor nonterminals.
All factor $h$-values can be cached with $O(n_K^2)$ multiplications in the
fixed monoid by extending factors one symbol at a time.

Rule~\textup{(B)} has at most $O(n_K^3)$ candidates: choose an observed
factor occurrence and a split point.  To enumerate Rule~\textup{(U)}, group
observed factor occurrences by their surrounding context $(u,v)$ and
$h$-value.  The total number of bucket entries is $O(n_K^2)$, while a fixed
context contains at most one factor from each sample word and therefore each
bucket has size at most $|K|\le n_K$.  Thus the total number of ordered pairs
emitted from all buckets is at most
\[
 \sum_B |B|^2
 \le n_K\sum_B |B|
 =O(n_K^3).
\]
Lexical and start rules contribute only lower-order terms.  With canonical
identifiers for equal factors and context components, all equality and type
tests after preprocessing are constant-time for fixed $h$.  Even if factor
strings are written literally in productions, each production has encoding
length $O(n_K)$, so the grammar can be checked and written explicitly in
$O(n_K^4)$ time.
\end{proof}

Each update of the conservative learner $\mathcal A_h$ is also polynomial
in the encoded data seen so far: CFG membership is polynomial in the grammar
and input sizes, and both the hypothesis size and any triggered rebuild are
polynomial by Theorem~\ref{thm:poly-build}.

Together with Lemma~\ref{lem:sample-consistency},
Theorem~\ref{thm:reconstruction-fixed-h} and Corollary~\ref{cor:ilt}, this
proves Theorem~\ref{thm:main}(i)--(iii).  The quantitative bounds in
items~(iv)--(v) are proved in Sections~\ref{sec:fixed-window-thick}--\ref{sec:linear}.
Grammar size alone cannot bound witness lengths: the standard
doubling grammar $S_0\to D_n$, $D_0\to a$,
$D_i\to D_{i-1}D_{i-1}$ has size $O(n)$ but generates the
singleton $\{a^{2^n}\}$, substitutable for every $h$;
compare~\cite[Example~2]{clark-eyraud2007} and
\cite[Example~2.3]{eyraud-heinz-yoshinaka2016}.

\section{Typed Thickness for General Fixed Typings}
\label{sec:typed-thickness}

For an arbitrary fixed $h$, the finite-witness construction admits a direct
quantitative bound once the cost of yield typing is measured explicitly.  Let
$G=(V,\Sigma,P,S_0)$ be a reduced SSBNF grammar and let $\widetilde G$ be its typed refinement.  Put
$N:=|V\setminus\{S_0\}|$ and
$N_t:=|\widetilde V\setminus\{\widetilde S\}|$
and define the \emph{typed thickness} of this presentation by
\[
 \tau_h^{\mathrm{typ}}(G):=
 \max\Bigl(\{\,|\omega(X)|:
 X\in\widetilde V\setminus\{\widetilde S\}\,\}\cup\{1\}\Bigr).
\]
Thus $\tau_h^{\mathrm{typ}}(G)$ measures shortest terminal yields after the
full yield-type split and trimming.  It is a presentation-dependent
parameter, not an additional assumption on the learner.

\begin{lemma}[typed-thickness witness bound]
\label{lem:typed-thickness-bound}
Let $G$ and $\widetilde G$ be as above and write
$\tau:=\tau_h^{\mathrm{typ}}(G)$.  Every non-start typed symbol $X$ has
$|u_X|+|v_X|\le (N_t-1)\tau$,
and every word $w\in\WitnessSet(\widetilde G)$ satisfies
$|w|\le (N_t+1)\tau$.
If $P_{\mathrm{ter}}$ and $P_{\mathrm{bin}}$ are respectively the terminal
and binary non-start productions of $G$, then
\[
 \|\WitnessSet(\widetilde G)\|
 \le
 \bigl(N|M|+|P_{\mathrm{ter}}|
       +|P_{\mathrm{bin}}||M|^2+1\bigr)
 \bigl((N|M|+1)\tau+1\bigr).
\]
In particular, for fixed $h$ this bound is polynomial in $|G|$ and
$\tau_h^{\mathrm{typ}}(G)$.
\end{lemma}

\begin{proof}
Form the directed dependency graph whose vertices are the non-start symbols of
$\widetilde G$ and whose edges join a parent to each child occurring in a
binary production.  Since $\widetilde G$ is reduced, every non-start symbol
$X$ is reachable.  Hence some start child $A$ has a path to $X$ in this
graph.  Choose a shortest such path.  It repeats no typed symbol, so it has at
most $N_t-1$ edges.

For each edge of the path, choose a witnessing binary production and expand
the child not lying on the path to its canonical yield.  Starting from
$\widetilde S\to A$, these expansions produce a terminal reaching context of
$X$.  Each off-path canonical yield has length at most $\tau$, so the total
context length is at most $(N_t-1)\tau$.  The canonical context
$(u_X,v_X)$ is no longer than this one.

A word in the first family adds one canonical yield to the canonical context;
the second family adds one terminal, and the third adds two canonical yields.
Since $\tau\ge1$, every word in~\eqref{eq:W} therefore has length at most
$(N_t+1)\tau$.  Finally, $N_t\le N|M|$; each source terminal production has
at most one typed terminal copy, and each source binary production has at
most $|M|^2$ typed copies.  Thus the displayed first factor bounds the number
of witnesses, including the possible witness $\lambda$, while the second
factor bounds $|w|+1$ for each witness.  Multiplying the two bounds proves the
claim.
\end{proof}

\begin{corollary}[characteristic data from typed thickness]
\label{cor:typed-thickness-data}
Fix $h$.  If $G$ is a reduced SSBNF grammar for a nonempty
$L(G)\in\Ccf{h}$, then $\WitnessSet(\widetilde G)$ is a characteristic
sample for $\mathcal B_h$ whose encoded size is polynomial in
$|G|$ and $\tau_h^{\mathrm{typ}}(G)$.
\end{corollary}

\begin{proof}
Lemma~\ref{lem:typed-thickness-bound} gives the quantitative bound, and
Theorem~\ref{thm:reconstruction-fixed-h} gives the characteristic-sample property.
\end{proof}

The following example shows a limitation of this
typed-refinement bound: typed thickness can grow exponentially
relative to source thickness even for a fixed $h$.

\begin{proposition}[exponential source-to-typed thickness gap]
\label{prop:typed-thickness-gap}
There is a fixed two-element typing $h_{\mathtt c}$ and a family of reduced
SSBNF grammars $G_n^{\mathtt c}$ of size $O(n)$ and ordinary thickness $1$
such that
$\tau_{h_{\mathtt c}}^{\mathrm{typ}}(G_n^{\mathtt c})\ge2^n$.
Consequently typed thickness cannot in general be bounded by a polynomial in
source grammar size and ordinary thickness, even for one fixed typing.
\end{proposition}
\begin{proof}

Let $M_{\mathtt c}=\{1,\zeta\}$ be the two-element monoid with
zero $\zeta$, with $h_{\mathtt c}(\mathtt a)=1$ and
$h_{\mathtt c}(\mathtt c)=\zeta$.  For $n\ge1$, define
$G_n^{\mathtt c}$ by
\[
\begin{aligned}
 S_0&\to U, &
 U&\to \mathtt a\mid\mathtt c\mid UU\mid E_nD, &
 D&\to\mathtt c,\\
 E_0&\to\mathtt a, &
 E_i&\to E_{i-1}E_{i-1}\mid\mathtt c
 \quad(1\le i\le n).
\end{aligned}
\]
These are reduced SSBNF grammars of size $O(n)$ with
$L(G_n^{\mathtt c})=\{\mathtt a,\mathtt c\}^+
\in\Ccf{h_{\mathtt c}}$.

Every nonterminal has a terminal yield of length $1$:
$S_0,U,E_0$ derive $\mathtt a$, and
$D,E_i$ ($i\ge1$) derive $\mathtt c$.
Hence the ordinary thickness is $\tau_{G_n^{\mathtt c}}=1$.

In contrast, after refinement to type $1$,
$E_i\to\mathtt c$ is unavailable because
$h_{\mathtt c}(\mathtt c)=\zeta\ne1$.
Since $1\cdot1$ is the only product equal to $1$,
induction gives $L((E_i)_1)=\{\mathtt a^{2^i}\}$.
Moreover, $\widetilde S\to U_\zeta\to(E_n)_1D_\zeta$
shows that $(E_n)_1$ survives trimming.
Thus $\tau_{h_{\mathtt c}}^{\mathrm{typ}}(G_n^{\mathtt c})
\ge2^n$.

\end{proof}

The target language $\{\mathtt a,\mathtt c\}^+$ in this family does not
depend on $n$.  Thus the proposition is not an algorithm-independent lower
bound on characteristic data.  Its role is only to show that ordinary source
thickness does not control the typed-refinement parameter appearing in
Corollary~\ref{cor:typed-thickness-data}; accordingly, the general fixed-$h$
bound is stated in terms of typed thickness.

\section{Fixed-Window and Locally Trivial Thick Data}
\label{sec:fixed-window-thick}

For fixed windows, preserving a bounded prefix and suffix controls
typed thickness by source thickness; this does not hold for arbitrary
$h$ (Proposition~\ref{prop:typed-thickness-gap}).
We obtain polynomial batch characteristic-data bounds in the sense of
identification in polynomial time and thick data (IPTtD) used by
Yoshinaka~\cite{yoshinaka2008} and Eyraud--Heinz--Yoshinaka
\cite[Definition~2.12]{eyraud-heinz-yoshinaka2016},
which adds grammar thickness as a parameter for characteristic-data
size in the set-based criterion.  We do not assert the same normal
form or polynomial degree.

Fix $k,\ell\ge0$, and let $G=(V,\Sigma,P,S_0)$ be a reduced SSBNF grammar.  Write
\[
 N:=|V\setminus\{S_0\}|,\qquad
 \theta_{k,\ell}(G):=
 \begin{cases}
   \tau_G, & k+\ell=0,\\[1mm]
   (k+\ell)+(2(k+\ell)-1)N\tau_G, & k+\ell>0,
 \end{cases}
\]
Let $\widetilde G$ be the reduced $h_{k,\ell}$-typed refinement of $G$.

\begin{lemma}[bounded fixed-window typed yields]
\label{lem:window-typed-yield}
Every non-start typed symbol $A_\mu$ of $\widetilde G$ satisfies
$|\omega(A_\mu)|\le\theta_{k,\ell}(G)$.
\end{lemma}

\begin{proof}
See Appendix~\ref{app:fixed-window-bounds}.
\end{proof}

\begin{theorem}[polynomial thick data for fixed windows]
\label{thm:window-thick}
For every fixed pair $(k,\ell)$ and ambient alphabet $\Sigma$, there is a
polynomial $q_{k,\ell}$ such that every reduced SSBNF grammar $G$ with
$L(G)\in\Ccf{h_{k,\ell}}$ satisfies
$\|\WitnessSet(\widetilde G)\| \le q_{k,\ell}(|G|,\tau_G)$,
and $\WitnessSet(\widetilde G)$ is a characteristic sample for
$\mathcal B_{h_{k,\ell}}$.
\end{theorem}

\begin{proof}
Lemma~\ref{lem:window-typed-yield} gives
$\tau_{h_{k,\ell}}^{\mathrm{typ}}(G)\le\theta_{k,\ell}(G)$, and for
fixed $k,\ell,\Sigma$ this is polynomial in $|G|$ and $\tau_G$.
Corollary~\ref{cor:typed-thickness-data} now gives the claim.
\end{proof}

The theorem is stated for reduced SSBNF representations.  The following
normalization transfers it to arbitrary reduced CFGs.

\begin{proposition}[polynomial thickness-preserving SSBNF normalization]
\label{prop:thick-ssbnf-normal}
There are fixed polynomials $q_{\mathrm{size}}$ and $q_{\mathrm{thick}}$ such that every reduced CFG
$G_*$ generating a nonempty language can be converted in polynomial time into
an equivalent reduced SSBNF grammar $G$ satisfying
$|G|\le q_{\mathrm{size}}(|G_*|)$ and
$\tau_G\le q_{\mathrm{thick}}(|G_*|,\tau_{G_*}+1)$.
\end{proposition}

The proof is in Appendix~\ref{app:thick-ssbnf}; terminal isolation and
binarization precede non-start $\lambda$-elimination so that grammar size and
shortest nonempty yields remain polynomially bounded.

\begin{corollary}[ordinary-thickness bound for locally trivial typings]
\label{cor:li-thickness}
Fix $h:\Sigma^*\to M$ and suppose that
$h(\Sigma^+)$ is locally trivial.  Then every nonempty
$L\in\Ccf{h}$ represented by a reduced CFG $G_*$ has characteristic data for
$\mathcal B_h$ whose encoded size is polynomial in $|G_*|$ and $\tau_{G_*}$.
\end{corollary}

\begin{proof}
Put $n:=|h(\Sigma^+)|+1$.  By Proposition~\ref{prop:li-window},
$\ker h_{n,n}|_{\Sigma^+}\subseteq\ker h|_{\Sigma^+}$.
For a reduced SSBNF grammar $G$, let $A_\mu$ be any symbol retained in the
reduced $h$-refinement.  Choose a successful derivation tree passing through
an occurrence of $A_\mu$ and annotate the same source tree by
$h_{n,n}$-yield types.  The corresponding copy $A_\nu$ lies on a successful
tree, hence is both reachable and productive and survives trimming.  By Proposition~\ref{prop:typed-core}, the yield languages of these symbols
are $L_G(A)\cap h_{n,n}^{-1}(\{\nu\})$ and
$L_G(A)\cap h^{-1}(\{\mu\})$, respectively.  The kernel refinement
$\ker h_{n,n}|_{\Sigma^+}\subseteq\ker h|_{\Sigma^+}$ makes the first
language a subset of the second.  Consequently
\[
 \tau_h^{\mathrm{typ}}(G)
 \le \tau_{h_{n,n}}^{\mathrm{typ}}(G)
 \le 2n+(4n-1)|V\setminus\{S_0\}|\tau_G
\]
by Lemma~\ref{lem:window-typed-yield}.  Since $h$ is fixed, $n$ is a
constant, and Corollary~\ref{cor:typed-thickness-data} gives a polynomial
characteristic-data bound for $G$.  Applying
Proposition~\ref{prop:thick-ssbnf-normal} transfers the bound to the
original reduced CFG $G_*$.
\end{proof}

The empty target is handled separately by the empty sample.

Theorem~\ref{thm:window-thick} and
Proposition~\ref{prop:thick-ssbnf-normal} also give the
fixed-window bound directly.  Polynomials may depend on the
fixed typing and alphabet; no uniform bound for variable
$k,\ell$, converse for non-locally-trivial typings, or
incremental polynomial-data guarantee is asserted.

\section{The Linear Subclass}
\label{sec:linear}

Linear languages are exactly the languages accepted by one-turn pushdown
automata~\cite{ginsburg-spanier1966,autebert1997}.  The normalization below
gives the grammatical counterpart: every successful derivation has one
continuing nonterminal path, which we call its \emph{spine}.

\begin{proposition}[linear-spine SSBNF normalization]\label{prop:linear-normal}
There is a fixed polynomial $q_{\mathrm{lin}}$ such that every linear CFG
$G$ generating a nonempty language can be converted in polynomial time into
an equivalent reduced SSBNF grammar $G_0=(V_0,\Sigma,P_0,S_0)$ with
$|G_0|\le q_{\mathrm{lin}}(|G|)$ in which every binary production has
exactly one fresh terminal-wrapper child $W_a$ with $W_a\to a$ and exactly
one other nonterminal child.
\end{proposition}

The construction is given in Appendix~\ref{app:linear-normalization}.  Let
$\widetilde G_0$ be its typed refinement and put
$N_t:=|\widetilde V_0\setminus\{\widetilde S\}|$.  Since each underlying
non-start symbol has at most $|M|$ typed copies,
$N_t\le|V_0\setminus\{S_0\}|\,|M|$.

\begin{lemma}[short canonical yields]\label{lem:linear-short}
Every typed non-start symbol $X$ of $\widetilde G_0$ satisfies
$|\omega(X)|\le N_t$.
Hence $\tau_h^{\mathrm{typ}}(G_0)\le N_t$.
\end{lemma}

The proof is given in Appendix~\ref{app:linear-short-proof}.

\begin{theorem}[polynomial characteristic data for linear targets]
\label{thm:linear-poly}
For fixed $h$, every $L\in\Clin{h}$ represented by a linear CFG $G$ has
a characteristic sample for $\mathcal B_h$ whose $\|\cdot\|$-size is
polynomial in $|G|$.
\end{theorem}

\begin{proof}
For $L=\emptyset$ and $L=\{\lambda\}$ use the samples from
Section~\ref{sec:typing}.  Otherwise take $G_0$ from
Proposition~\ref{prop:linear-normal}.  Lemma~\ref{lem:linear-short} gives
$\tau_h^{\mathrm{typ}}(G_0)\le N_t\le
|V_0\setminus\{S_0\}|\,|M|$, and
Corollary~\ref{cor:typed-thickness-data} yields a characteristic sample
polynomial in $|G_0|\le q_{\mathrm{lin}}(|G|)$.
\end{proof}

Together with Theorem~\ref{thm:poly-build}, this gives
polynomial-time batch reconstruction and polynomial characteristic data
relative to a chosen linear CFG for targets in $\Clin{h}$.
This does not contradict de~la~Higuera's negative result for
the unrestricted class of linear languages
\cite[Theorem~2]{higuera1997}: our theorem concerns only the
fixed-$h$ substitutable subclass.  Moreover, no incremental
polynomial-data guarantee for $\mathcal A_h$ is asserted.

\paragraph{A nonregular linear \texorpdfstring{fixed-$h$}{fixed-h} example.}

The preceding theorem is not confined to regular targets, nor to the
Clark--Eyraud or fixed-window classes.  Define
\[
\begin{split}
L_{\pm,e}:={}&
\{a^{2m}cb^{2m}:m\ge0\}
\cup
\{a^{2m+1}db^{2m+1}:m\ge0\}\\
&\cup
\{a^neb^n:n\ge0\}.
\end{split}
\]
It is generated by the linear grammar
\[
\begin{array}{lll}
S\to U_{\mathrm{ev}}\mid U_e, &
U_{\mathrm{ev}}\to c\mid aU_{\mathrm{odd}}b, &
U_{\mathrm{odd}}\to d\mid aU_{\mathrm{ev}}b,\\[1mm]
&&U_e\to e\mid aU_eb.
\end{array}
\]

The subscripts $\mathrm{ev}$ and $\mathrm{odd}$ track the parity
of enclosing $a,b$-pairs from the start, not the length parity
of the corresponding nonterminal's own yield.


Call $c,d,e$ the \emph{center symbols}.  Let
$M_{\pm,e}:=\{1,\kappa_{cd},\kappa_e,0\}$ be the four-element monoid in
which $1$ is the identity and $\mu\nu=0$ for all
$\mu,\nu\in\{\kappa_{cd},\kappa_e,0\}$.  Thus $M_{\pm,e}$ is obtained
by adjoining an identity to a three-element null semigroup.  Define the
homomorphism $h_{\pm,e}:\{a,b,c,d,e\}^*\to M_{\pm,e}$ by
\begin{equation}\label{eq:lpm-h}
 h_{\pm,e}(a)=h_{\pm,e}(b)=1,\qquad
 h_{\pm,e}(c)=h_{\pm,e}(d)=\kappa_{cd},\qquad
 h_{\pm,e}(e)=\kappa_e.
\end{equation}
Hence $h_{\pm,e}(w)=1$ if $w$ contains no center symbol, $h_{\pm,e}(w)=\kappa_{cd}$ if it
contains exactly one center symbol and that symbol is $c$ or $d$,
$h_{\pm,e}(w)=\kappa_e$ if that unique center is $e$, and $h_{\pm,e}(w)=0$ if it contains at
least two center symbols.
The choice $h_{\pm,e}(c)=h_{\pm,e}(d)$ leaves the $c/d$
parity distinction to distributional behavior; the separate
$e$-type gives the fixed-window collision below.

Put $\operatorname{par}(c):=0$ and $\operatorname{par}(d):=1$.  Then, for
$m,n\ge0$ and $\xi\in\{c,d,e\}$,
\begin{equation}\label{eq:lpm-membership}
a^m\xi b^n\in L_{\pm,e}
\iff
m=n
\ \text{and}\
\bigl[
 \xi=e
 \ \text{or}\
 (\xi\in\{c,d\}
  \ \text{and}\
  m\equiv\operatorname{par}(\xi)\pmod2)
\bigr].
\end{equation}

\begin{proposition}\label{prop:linear-separator-example}
The language $L_{\pm,e}$ is linear and nonregular and is
$\sim_{h_{\pm,e}}$-substitutable for the homomorphism in~\eqref{eq:lpm-h}.  On the
other hand, it is not Clark--Eyraud substitutable and is not
$(k,\ell)$-substitutable for any fixed $k,\ell\ge0$.
\end{proposition}

\begin{proof}
Linearity is immediate from the displayed grammar, while
$L_{\pm,e}\cap a^*eb^*=\{a^neb^n:n\ge0\}$ proves nonregularity.  Write
$L:=L_{\pm,e}$.

For $\sim_{h_{\pm,e}}$-substitutability, let equal-type nonempty factors
$x,y$ share a context in $L$.  No factor occurring in a word of $L$ can have
type $0$.  If their common type is $1$, then
$x,y\in a^+\cup b^+$; the common context forces
$|x|_a-|x|_b=|y|_a-|y|_b$, which determines a word in
$a^+\cup b^+$ uniquely, so $x=y$.

Otherwise write $x=a^i\xi b^j$ and $y=a^{i'}\xi'b^{j'}$ with one center
symbol each.  By~\eqref{eq:lpm-membership}, every context of such a factor
has the form $(a^m,b^n)$, and
\[
\begin{array}{ll}
(a^m,b^n)\in\D_L(a^ieb^j)
 &\Longleftrightarrow n-m=i-j,\\
(a^m,b^n)\in\D_L(a^i\xi b^j),\ \xi\in\{c,d\}
 &\Longleftrightarrow
 n-m=i-j\ \text{and}\
 m\equiv\operatorname{par}(\xi)-i\pmod2.
\end{array}
\]
Thus an $e$-factor's distribution is determined by $i-j$, and a $c/d$-factor's
by $(i-j,\operatorname{par}(\xi)-i\bmod2)$.  A shared context makes the corresponding invariant equal for $x$ and $y$;
equality of $h_{\pm,e}$-types ensures that either both are $e$-factors or
both are $c/d$-factors.  Hence $\D_L(x)=\D_L(y)$.

The language is not Clark--Eyraud substitutable because $c,e\in L$ share the
empty context but $(a,b)\in\D_L(e)\setminus\D_L(c)$.  Finally fix
$k,\ell$ and put $m_0:=\max\{1,k+\ell\}$,
$w_e:=a^keb^\ell$, $w_c:=a^kcb^\ell$.  These factors have the same
length-$k$ prefix and length-$\ell$ suffix and share
$(a^{2m_0-k},b^{2m_0-\ell})$, whereas
$(a^{2m_0+1-k},b^{2m_0+1-\ell})
 \in\D_L(w_e)\setminus\D_L(w_c)$.
The \hyperref[crit:kl]{fixed-window counterexample criterion} therefore gives
$L_{\pm,e}\notin\mathsf{KL}$.
\end{proof}

This example also arises from regular filtering.  Put
\[
 L_{\mathrm{all}}=\{a^n\xi b^n:n\ge0,\ \xi\in\{c,d,e\}\},
 \quad Q=(aa)^*c(bb)^*\cup a(aa)^*d\,b(bb)^*\cup a^*eb^*.
\]
The linear $L_{\mathrm{all}}$ is Clark--Eyraud substitutable:
a common context fixes a center-free factor, and one-center
factors' distributions depend only on $i-j$.  As
$L_{\pm,e}=L_{\mathrm{all}}\cap Q$, this exemplifies
Proposition~\ref{prop:finite-info-closure}(ii) and closure of
linear languages under regular intersection:
\[
 L\in\Clin{h},\quad Q=g^{-1}(F)
 \quad\Longrightarrow\quad L\cap Q\in\Clin{h\times g}.
\]
The direct four-element typing~\eqref{eq:lpm-h} is smaller
than the general product typing.

Thus Theorem~\ref{thm:linear-poly} applies to a nonregular linear language
outside $\mathsf{KL}$.  The displayed grammar is even-linear, so no
separation from Takada's control-set framework is claimed.

\paragraph{Finite-state control induced by yield typing.}
The following straightforward product construction gives a limited
connection with Takada's control-set viewpoint. Give each production of a linear grammar $G=(V,\Sigma,P,S)$ a distinct label
from $\Pi$, and write $A\Rightarrow_G^\alpha w$ when
$\alpha\in\Pi^*$ is the production-label word of a derivation.  For
$F\subseteq M$, let
\[
 C_{h,F}(G):=
 \{\alpha\in\Pi^*:S\Rightarrow_G^\alpha w
   \text{ for some }w\in\Sigma^*,\ h(w)\in F\}.
\]
This is a regular control set.  A finite automaton uses states
$V\times M\times M$ (and one accepting state), starts in $(S,1,1)$, and for
a labeled rule $A\to uBv$ makes the transition
$(A,\mu_L,\mu_R)\mapsto
(B,\mu_L h(u),h(v)\mu_R)$.  A terminal rule $A\to w$ is accepted
exactly when $\mu_L h(w)\mu_R\in F$.  The invariant is that a sentential
form $uAv$ is represented by $(A,h(u),h(v))$, so the accepted label words are
exactly those of successful derivations whose yields have $h$-type in $F$.
Thus regular control by $C_{h,F}(G)$ generates precisely
$L(G)\cap h^{-1}(F)$.  Starting from any nonterminal $A$ with
$F=\{\mu\}$ gives the same interpretation for a typed copy $A_\mu$.
This is a finite-state bridge to Takada's regular control
languages~\cite{takada1995}; it does not claim that arbitrary Takada control
sets arise from terminal yield typings.

\section{Expressiveness and Boundaries}
\label{sec:boundaries}

We now locate the framework between the fixed-window hierarchy and broader
context-free families.

\subsection{A nonlinear target inside \texorpdfstring{$\mathsf{RS}$}{RS}}

Let $\Delta:=\{a^nb^n:n\ge0\}$ and consider $\Delta^*$.  For
$z\in\{a,b\}^*$ put $\beta(z):=|z|_a-|z|_b$ and define
\[
 \underline\beta(z):=
 \min\{\beta(p):p\text{ is a prefix of }z\},
\]
where the empty prefix is included.  For an occurrence of $ba$, define its
height to be the height of the prefix ending at that $b$.
We use:
\textup{($\star$)} $w\in\Delta^*$ iff
$\underline\beta(w)\ge0$, $\beta(w)=0$, and every occurrence of $ba$
has height $0$.  The forward implication follows from the factorization into
$a^nb^n$ factors; conversely, cutting at successive zero-height positions
gives balanced $a^*b^*$ segments.

Let $\operatorname{fst}(w),\operatorname{lst}(w)\in\{\bot,a,b\}$ be
the first and last symbols, with $\bot$ on $\lambda$, and let
$\operatorname{ba}(w)$ record whether $w$ contains $ba$.  Put
\[
h_\star(w):=
(\operatorname{fst}(w),\operatorname{lst}(w),\operatorname{ba}(w)).
\]
The summary of a concatenation is determined by the summaries of its factors:
the endpoints concatenate in the evident way and a new $ba$ can arise only
across a boundary $b|a$.  Hence the finite image
$M_\star:=h_\star(\{a,b\}^*)$ carries the product
$h_\star(x)h_\star(y):=h_\star(xy)$, making $h_\star$ a finite-monoid
homomorphism.

For nonempty $x$, let $\operatorname{Occ}_{ba}(x)$ denote its
internal occurrences of $ba$.  For each $o\in\operatorname{Occ}_{ba}(x)$,
let $\eta_x(o)$ be the height, relative to the beginning of $x$,
of the prefix ending at the $b$ in this occurrence.  Define
\[
 \mathcal Q(x):=
 \{q\in\mathbb Z:
 q+\underline\beta(x)\ge0\ \text{and}\
 q+\eta_x(o)=0\text{ for all }o\in\operatorname{Occ}_{ba}(x)\}.
\]

\begin{proposition}[a nonlinear fixed-$h_\star$ target]
\label{prop:nonlinear-rs-example}
The language $\Delta^*$ is deterministic context-free, nonregular, and
nonlinear, and it is $\sim_{h_\star}$-substitutable.  Moreover it is not
$(k,\ell)$-substitutable for any fixed $k,\ell\ge0$.  Hence
$\Delta^*\in\mathsf{RS}_{h_\star}(\{a,b\})\cap\mathrm{CFL}$ and
$\Delta^*\in\mathsf{RS}(\{a,b\})\setminus\mathsf{KL}(\{a,b\})$.
\end{proposition}

\begin{proof}
The grammar $S\to TS\mid\lambda$, $T\to aTb\mid\lambda$ generates
$\Delta^*$.  It is deterministic context-free: a DPDA pushes while reading
the $a$-part of a factor $a^nb^n$, pops during its $b$-part, and
permits a new such factor only after the stack returns to its bottom marker.  Nonregularity
follows from
$\Delta^*\cap a^*b^*=\Delta$, and nonlinearity from closure of linear
languages under regular intersection together with
$\Delta^*\cap a^*b^*a^*b^*=\Delta\Delta$.
The nonlinearity of $\Delta\Delta$ is explicitly recorded by
Autebert--Berstel--Boasson~\cite[Section~6.1, before Theorem~6.1]{autebert1997}.

For substitutability, let nonempty $x,y$ have equal $h_\star$-type and a
common context $(u,v)$.  Since $uxv,uyv\in\Delta^*$, total balance gives
$\beta(x)=\beta(y)$.  Fix an arbitrary context $(u',v')$.  By
\textup{($\star$)}, the conditions contributed while traversing the middle
factor $x$ are exactly that the height never becomes negative inside $x$ and
that every internal occurrence of $ba$ in $x$ is met at global height $0$;
these are equivalent to $\beta(u')\in\mathcal Q(x)$.  A possible new
occurrence of $ba$ across either boundary $u'|x$ or $x|v'$ depends only on
the first and last symbols of $x$, while the height at which $v'$ is entered
depends only on $\beta(u')+\beta(x)$.  Thus, once $u'$ and $v'$ are fixed,
all remaining membership conditions depend on $x$ only through
$h_\star(x)$, $\beta(x)$, and $\mathcal Q(x)$.  It therefore suffices to
show $\mathcal Q(x)=\mathcal Q(y)$.

If both contain $ba$, any internal $ba$ forces the unique possible entry
height, and the common context supplies it; hence
$\mathcal Q(x)=\{\beta(u)\}=\mathcal Q(y)$.  If neither contains
$ba$, then $x,y\in a^*b^*$ and
$\mathcal Q(z)=\{q:q\ge-\min\{0,\beta(z)\}\}$, so equality follows
from $\beta(x)=\beta(y)$.  Thus
$\D_{\Delta^*}(x)=\D_{\Delta^*}(y)$.

For the fixed-window separation, fix $k,\ell$ and let
$m:=\max\{1,k+\ell\}$,
$x:=a^mb^m$, and $y:=a^mb^ma^mb^m$.  They have the same length-$k$ prefix
and length-$\ell$ suffix and share the empty context.  Moreover
$axb\in\Delta^*$, whereas the central $ba$ in $ayb$ has height $1$, so
$ayb\notin\Delta^*$.  The fixed-window counterexample criterion therefore
excludes every $(k,\ell)$.
\end{proof}


\subsection{A capped regular counter family}

Yoshinaka's regular language
$L_0=ae^*ce^*a\cup ae^*de^*a\cup be^*ce^*b$
already gives a fixed-window separation
\cite[Proposition~1]{yoshinaka2008}; capped counters contrast
with the uncapped obstruction below.

Let $\Sigma_c:=\{\uparrow,\downarrow,\#\}$.  For
$w=\delta_1\cdots\delta_n\in\Sigma_c^*$, let $\mathrm{ht}_i(w)$
be the height after $i$ letters, with $\mathrm{ht}_0(w)=0$, changing
by $+1$ on $\uparrow$, $-1$ on $\downarrow$, and resetting to $0$ on $\#$.
For $\rho\ge1$, put
\[
\mathrm{CTR}_\rho:=
\{w\in\Sigma_c^*:0\le\mathrm{ht}_i(w)\le\rho
\text{ for all }0\le i\le|w|\}.
\]

For every $\rho\ge1$, $\mathrm{CTR}_\rho$ is regular: a deterministic
finite automaton records
heights $0,\ldots,\rho$, uses one sink outside the range, and resets to $0$
on $\#$.  Hence $\mathrm{CTR}_\rho\in\mathsf{RS}(\Sigma_c)$ by
Proposition~\ref{prop:regular-auto}.

\begin{theorem}\label{thm:ctr-non-kl}
For every $\rho\ge2$, $\mathrm{CTR}_\rho$ is not
$(k,\ell)$-substitutable for any $k,\ell\ge0$.
\end{theorem}
\begin{proof}
Fix $k,\ell$ and put
$x:=(\uparrow\downarrow)^k\uparrow^{\rho-1}\#^\ell$ and
$y:=(\uparrow\downarrow)^k\uparrow^\rho\#^\ell$.  They are long
enough for the fixed-window criterion, have the same length-$k$ prefix and
length-$\ell$ suffix, and both lie in $\mathrm{CTR}_\rho$.  Yet
$\uparrow x\in\mathrm{CTR}_\rho$ while
$\uparrow y\notin\mathrm{CTR}_\rho$, the latter reaching height
$\rho+1$.  Hence their distributions differ, and the criterion applies.
\end{proof}

The bound is sharp: $\mathrm{CTR}_1$ is $(1,1)$-substitutable,
since membership is determined by the initial letter and the
forbidden adjacent pairs $\uparrow\uparrow$,
$\downarrow\downarrow$, and $\#\downarrow$.
The DFA transition monoid gives a typing $h_\rho$
recognizing each regular $\mathrm{CTR}_\rho$; only $\rho\ge2$
lies outside the fixed-window hierarchy.

\subsection{Natural deterministic context-free languages outside \texorpdfstring{$\mathsf{RS}$}{RS}}\label{sec:outside-rs}

\begin{lemma}[finite-monoid obstruction]\label{lem:finite-monoid-obstruction}
Let $L\subseteq\Sigma^*$.  Suppose an infinite
$\Xi\subseteq\Sigma^+$ satisfies, for all distinct $x,y\in\Xi$,
$\D_L(x)\cap\D_L(y)\ne\emptyset$ and
$\D_L(x)\ne\D_L(y)$.  Then $L\notin\mathsf{RS}(\Sigma)$.
\end{lemma}
\begin{proof}
Any finite $h$ identifies two elements of $\Xi$; their common context and
unequal distributions contradict $\sim_h$-substitutability.
\end{proof}

Removing the upper cap gives the deterministic context-free language
\[
\mathrm{CTR}:=
\{w\in\Sigma_c^*: \mathrm{ht}_i(w)\ge0
\text{ for all }0\le i\le|w|\},
\]
accepted by the evident deterministic pushdown counter with reset.

\begin{corollary}[uncapped counter obstruction]
$\mathrm{CTR}\notin\mathsf{RS}(\Sigma_c)$.
\end{corollary}
\begin{proof}
Use $\Xi=\{\uparrow^i:i\ge1\}$.  For $i<j$,
$(\lambda,\downarrow^i)$ is common, while
$(\lambda,\downarrow^j)$ is a context only of $\uparrow^j$.
\end{proof}

For the one-bracket Dyck language
\[
 D_1:=\{w\in\{a,b\}^*: |w'|_a\ge|w'|_b
 \text{ for every prefix }w',\ |w|_a=|w|_b\},
\]
a standard deterministic context-free language
\cite{ginsburg-greibach1966}, our notation denotes the whole Dyck language
(called $D_1^*$ in the convention of
Autebert--Berstel--Boasson~\cite{autebert1997}).

\begin{corollary}\label{cor:dyck-not-rs}
$D_1\notin\mathsf{RS}(\{a,b\})$.
\end{corollary}
\begin{proof}
Use $\Xi=\{b^ia^i:i\ge1\}$.  For $i<j$,
$(a^j,b^j)$ is common, whereas $(a^i,b^i)$ is a context only of
$b^ia^i$; for $b^ja^j$ the prefix $a^ib^j$ underflows.
\end{proof}

Thus unbounded stack behavior alone is not the obstruction: the nonregular
one-turn language $L_{\pm,e}$ lies in $\mathsf{RS}$, but an infinite family
of mutually comparable factors with distinct contextual behavior cannot be
separated by any finite $h$.

\begin{lemma}[fixed-word quotient]\label{lem:rs-fixed-quotient}
If $L$ is $\sim_h$-substitutable and $z\in\Sigma^*$, then
$L/z:=\{w:wz\in L\}$ is $\sim_h$-substitutable.
\end{lemma}
\begin{proof}
A common context $(u,v)$ in $L/z$ becomes $(u,vz)$ in $L$; equality of
distributions in $L$ then transfers every context $(s,t)$ by appending $z$.
\end{proof}

The language generated by $S\to aSS\mid b$ is the
\emph{{\fontencoding{T1}\selectfont\L{}}ukasiewicz language}.
Yoshinaka~\cite[Section~3]{yoshinaka2008} already observed that this language
is not $(k,\ell)$-substitutable for any fixed $k,\ell$.  In the notation used
here, $L_{\mathrm{Luk}}=D_1b$~\cite{autebert1997}, and it is deterministic
context-free~\cite{ginsburg-greibach1966}.  The quotient argument gives the
stronger conclusion that it lies outside the present finite-typing union: if
$L_{\mathrm{Luk}}\in\mathsf{RS}(\{a,b\})$, then
Lemma~\ref{lem:rs-fixed-quotient} would put
$D_1=L_{\mathrm{Luk}}/b$ in $\mathsf{RS}$, contradicting
Corollary~\ref{cor:dyck-not-rs}.

\subsection{Comparison with Clark's congruential family}

Clark~\cite{clark2010congruential} allows a finite initial set
$I\subseteq V$; write $L_G[I]:=\bigcup_{A\in I}L_G(A)$.  Such a grammar is
\emph{congruential} if each $L_G(A)$ lies in one syntactic-congruence class
of $L_G[I]$.  Let $\mathsf{CONG}(\Gamma)$ be the languages over $\Gamma$
generated by such grammars.  Clark already places substitutable and
$(k,\ell)$-substitutable languages in this family; the following extends
the containment to every fixed-$h$ slice.

\begin{proposition}[comparison with Clark's congruential family]
\label{prop:clark-congruential-comparison}
For every finite alphabet $\Gamma$ and finite-monoid homomorphism
$h:\Gamma^*\to M$,
$\Ccf{h}\subseteq\mathsf{CONG}(\Gamma)$.  Consequently
$\mathsf{RS}(\Gamma)\cap\mathrm{CFL}
\subseteq\mathsf{CONG}(\Gamma)$, and for $\Gamma=\{a,b\}$ the
inclusion is proper.
\end{proposition}

\begin{proof}
Let nonempty $L\in\Ccf{h}$ and take its reduced SSBNF grammar and typed
refinement $\widetilde G$.  Every non-start typed symbol $X=A_\mu$ has a
terminal reaching context $(u,v)$, and every
$w\in L_{\widetilde G}(X)$ is nonempty with $h(w)=\mu$ by
Proposition~\ref{prop:typed-core}.  Hence any
$x,y\in L_{\widetilde G}(X)$ have the same type and the common context
$(u,v)$; fixed-$h$ substitutability gives $x\equiv_L y$.  Thus each typed
nonterminal generates strings from one syntactic-congruence class.

Delete the separated start symbol and use its typed successors as the finite
initial set; if $\lambda\in L$, add one fresh initial symbol generating only
$\lambda$.  Proposition~\ref{prop:typed-core} then gives a congruential
grammar for $L$.  The empty language is handled by an initial symbol with no
productions.  This proves the inclusion after taking the union over $h$.

For properness over $\{a,b\}$, use the Dyck example also discussed by
Clark~\cite{clark2010congruential}.  Every $x\in D_1$ satisfies
$uxv\in D_1$ iff $uv\in D_1$, so $x\equiv_{D_1}\lambda$.  Hence
$S\to aSbS\mid\lambda$ is congruential, while
Corollary~\ref{cor:dyck-not-rs} gives
$D_1\notin\mathsf{RS}(\{a,b\})$.
\end{proof}

\section{Conclusion}

For fixed $h$, we separate polynomial-time set-driven reconstruction
from conservative Gold identification.  Canonical witnesses yield
characteristic data bounded by typed thickness; ordinary thickness
suffices for locally trivial positive-image typings, while linear
targets admit a grammar-size bound.  The typed-thickness gap
limits this refinement analysis, not learning algorithms in general.

After finite refinement of the typing, product typing accommodates regular
filtering and arbitrary inverse homomorphisms; this is a compositionality
statement rather than closure of one fixed-$h$ slice.  The locally trivial
criterion recovers the fixed-window union.
Nonregular linear and nonlinear examples reach beyond fixed windows;
counter/Dyck obstructions and strict containment in Clark's
congruential family delimit the framework.  Open questions include
characterizing $\mathsf{RS}\cap\mathrm{CFL}$, the extent of
$\Clin{h}$ for nonregular targets, how to choose useful typings,
and direct pushdown characterizations.

\section*{Data availability}
No empirical data were used for the research described in this article.

\section*{Declaration of competing interest}
The author declares no known competing financial interests or personal
relationships that could have influenced the work reported in this paper.

\section*{Funding}
This research received no specific grant from any funding agency in the
public, commercial, or not-for-profit sectors.

\section*{Acknowledgments}
The author thanks Makoto Kanazawa and Ryo Yoshinaka for guidance,
encouragement, and help with relevant literature, and the two anonymous
reviewers for careful and constructive reports that substantially improved
the paper, in particular by identifying problems in the submitted learning
algorithm and by prompting the direct two-word witness in
Theorem~\ref{thm:ctr-non-kl}.  Any remaining errors are the author's
responsibility.

\appendix
\section{Fixed-window witness bounds and thickness-preserving normalization}
\label{app:fixed-window-bounds}

\subsection{Proof of Lemma~\ref{lem:window-typed-yield}}

\begin{proof}[Proof of Lemma~\ref{lem:window-typed-yield}]
When $k+\ell=0$, $h_{0,0}$ is constant on $\Sigma^+$, so a non-start typed
symbol $A_\mu$ of $\widetilde G$ has the same nonempty yields as its
underlying nonterminal $A$ and hence $|\omega(A_\mu)|\le\tau_G$.

Assume $k+\ell>0$.  If $\mu$ is a short-word
element of $M_{k,\ell}$, then every word of type $\mu$ is that very word, of
length $<k+\ell$, and the required bound is immediate.  It remains to consider a
long-word summary element $\mu=\langle p,q\rangle$.  Choose any derivation
tree of $A$ whose terminal yield $w$ has this type.  Mark the first $k$
terminal leaves and the last $\ell$ terminal leaves; there are exactly
$k+\ell$ marked leaves, and their labels are $p$ and $q$ in the required order.

Let $\mathcal U$ be the union of the root-to-marked-leaf paths.  Consider each
maximal chain of nonterminal nodes in $\mathcal U$ along which exactly one child
remains in $\mathcal U$.  If the same underlying nonterminal label occurs twice on
such a chain, replace the derivation rooted at the upper occurrence by the
subtree rooted at the lower occurrence.  This is a valid derivation from the
same nonterminal and deletes only sibling subtrees containing no marked leaf,
so all marked boundary terminals and their order are preserved.  Repeating
the shortcut leaves at most $N$ nonterminal nodes on each maximal one-child
chain.

Contract every maximal one-child chain.  The resulting rooted binary tree
connects the $k+\ell$ marked leaves and has at most $2(k+\ell)-1$ vertices,
so there are at most $2(k+\ell)-1$ maximal one-child chains (including a possible initial chain
above the first branching vertex).  An omitted sibling subtree can occur only
at a node on one of these chains; branching nodes have both children in
$\mathcal U$.  Hence there are at most $(2(k+\ell)-1)N$ omitted sibling subtrees.  Replace
each by a shortest terminal yield of its root, contributing at most $\tau_G$
terminals per subtree.  Since the marked leaves are exactly the first $k$
and last $\ell$ leaves, no unmarked sibling subtree lies before the marked
prefix (when $k>0$) or after the marked suffix (when $\ell>0$).
Every non-start yield in SSBNF is nonempty.  Thus shortcuts and sibling
replacements preserve the marked prefix and suffix, even though they may
alter the material between them; if $k=0$ or $\ell=0$, the corresponding
boundary condition is vacuous.  The resulting word $w'$ is still derived from $A$,
still has prefix $p$ and suffix $q$, and therefore has the same
$h_{k,\ell}$-type $\mu$, while
$|w'|\le (k+\ell)+(2(k+\ell)-1)N\tau_G$.
By the exact lifting equality in Proposition~\ref{prop:typed-core},
$w'\in L_G(A)\cap h_{k,\ell}^{-1}(\{\mu\})$ is a typed yield of the
retained $A_\mu$ in $\widetilde G$.  Hence the shortlex-minimal typed yield
$\omega(A_\mu)$ obeys the same length bound.
\end{proof}

\subsection{Thickness-preserving SSBNF normalization}
\label{app:thick-ssbnf}

\begin{proof}[Proof of Proposition~\ref{prop:thick-ssbnf-normal}]
Let $G_*=(V,\Sigma,P,S)$ be reduced and generate a nonempty language.  Put
$n:=|G_*|$ and $\bar\tau:=\tau_{G_*}+1$.  We describe a standard normalization in
an order that keeps both grammar size and shortest terminal yields polynomial.

First introduce a fresh start symbol $S_0$ with $S_0\to S$ and, if
$\lambda\in L(G_*)$, retain the possibility $S_0\to\lambda$ at the end.  Isolate every terminal occurring in a right-hand side of length at least two
by a fresh wrapper $W_a\to a$, and binarize every right-hand side of length
greater than two.
Call the resulting grammar $G_{\mathrm{bin}}$.  These operations have total size $O(n)$ up
to the chosen reasonable encoding.  Every original nonterminal keeps the same
terminal language.  Every fresh binarization symbol expands a contiguous block
of some original right-hand side, containing at most $O(n)$ original symbols;
replacing each productive nonterminal in that block by a shortest terminal
yield gives
$\tau_{G_{\mathrm{bin}}}\le c_1 n\bar\tau$
for a fixed constant $c_1$.

We next bound the shortest \emph{nonempty} yield of a nullable symbol of $G_{\mathrm{bin}}$.
Suppose $A\Rightarrow_{G_{\mathrm{bin}}}^*w$ for some nonempty terminal word $w$.  Choose one
terminal leaf in a derivation tree and follow the path from the root to that
leaf.  If the same nonterminal occurs twice on this path, replace the subtree
at the upper occurrence by the subtree at the lower occurrence.  Repeating
this shortcut yields a path with at most $|V_{G_{\mathrm{bin}}}|$ nonterminal occurrences.
At every binary node on the path, replace the off-path sibling by a shortest
terminal yield; a nullable sibling may contribute the empty word.  The chosen
terminal leaf contributes one symbol.  Hence
\[
 \min\{\,|w|>0:A\Rightarrow_{G_{\mathrm{bin}}}^*w\,\}
 \le 1+|V_{G_{\mathrm{bin}}}|\tau_{G_{\mathrm{bin}}}
 \le c_2 n^2\bar\tau
\]
for another fixed constant $c_2$.

Now eliminate all non-start $\lambda$-productions.  Because $G_{\mathrm{bin}}$ is binary, a
rule $A\to BC$ creates at most the variants $A\to BC$, $A\to B$, and
$A\to C$ according to which children are nullable; there is therefore no
exponential nullable-subset expansion.  The usual construction preserves, for
every nonterminal, its nonempty terminal language, so the preceding bound is
unchanged.  Next compute unit closure and eliminate non-start unit productions.
For each source nonterminal and each unit-reachable non-unit production, copy
that production to the source.  With $O(n)$ nonterminals and $O(n)$
non-unit productions before this step, this creates at most $O(n^2)$ rules.
Unit elimination preserves every nonterminal language and therefore does not
increase the thickness bound.  Finally trim useless symbols.

The resulting grammar is reduced and has separated start form; every non-start
rule is terminal or binary, and $S_0\to\lambda$ is the only possible
$\lambda$-rule.  Thus it is in SSBNF.  Its size is polynomial in $n$, and
its thickness is at most $c_2 n^2(\tau_{G_*}+1)$.  All transformations above are
computable in polynomial time, proving the proposition.
\end{proof}

\section{Proof of the linear-spine SSBNF normalization}
\label{app:linear-normalization}

The construction for Proposition~\ref{prop:linear-normal} combines standard
start-symbol separation, elimination of useless, nullable, and unit rules,
and binarization, with the cases written out to preserve the single-spine form
needed in Section~\ref{sec:linear}.

\begin{proof}[Proof of Proposition~\ref{prop:linear-normal}]
If $L(G)=\{\lambda\}$, take the reduced grammar with the single rule
$S_0\to\lambda$; it has no binary productions, so the required condition is
vacuous.  Hence assume below that $L(G)$ contains a nonempty word.

Let $S$ be the original start symbol.  Introduce a fresh $S_0$ that never
occurs on a right-hand side, with $S_0\to S$ and, if $\lambda\in L(G)$,
$S_0\to\lambda$.  Treat $S$ thereafter as a non-start symbol.  Remove useless
symbols, eliminate non-start $\lambda$-rules while preserving $\lambda$ only
at $S_0$, and then eliminate non-start unit rules.  Unit elimination is not
applied to the start rules, so every start rule surviving final trimming has
the required form $S_0\to A$ or $S_0\to\lambda$.  This order also removes unit rules created by $\lambda$-elimination.
Linearity implies that nullable elimination creates at most one additional
alternative per production and leaves at most one nonterminal on every
surviving right-hand side.  Unit closure copies at most one non-unit
production for each source/reachable-production pair, so these preprocessing
steps remain polynomial.

Binarize using fresh terminal wrappers $W_a\to a$, without
introducing unit rules.  For each rule
$A\to a_1\cdots a_m B b_1\cdots b_n$, unit elimination ensures
$m+n>0$.  Put $R_0:=B$ and, for $n>0$, form a right chain
\[
 R_j\to R_{j-1}W_{b_j}\qquad(1\le j\le n).
\]
All intermediate $R_j$ are fresh.  If $m=0$ let $R_n:=A$,
so the last displayed rule has left side $A$; otherwise $R_n$
is fresh for $n>0$, and $R_0=B$ when $n=0$.
For $m>0$ put $L_m:=R_n$, introduce fresh
$L_1,\ldots,L_{m-1}$, and use
\[
 A\to W_{a_1}L_1,\qquad
 L_i\to W_{a_{i+1}}L_{i+1}\quad(1\le i<m).
\]
Here $L_1=R_n$ when $m=1$.  Each rule is binary
with precisely one fresh terminal-wrapper child;
the original symbol $B$ is not a wrapper.

For a terminal-only rule $A\to a_1\cdots a_m$, keep $A\to a_1$ when
$m=1$.  For $m\ge2$, introduce fresh
$\Theta_2,\ldots,\Theta_m$ and use
\[
 A\to W_{a_1}\Theta_2,\qquad
 \Theta_i\to W_{a_i}\Theta_{i+1}\quad(2\le i\le m-1),\qquad
 \Theta_m\to a_m.
\]
(When $m=2$, the middle family is empty.)  Every new non-start rule is
therefore either terminal or binary with exactly one wrapper child and one
other nonterminal child.  By construction, a wrapper $W_a$ occurs only in the
wrapper-child position of a binary rule and never as its continuing
nonterminal child; the terminal-only construction ends instead in a fresh non-wrapper
nonterminal with a terminal rule.  Final trimming preserves this property,
the separated start form, and reducedness.  The wrapper and chain
factorizations are linear in the already polynomial total right-hand-side
length, so the whole construction has polynomial size and runs in polynomial
time.
\end{proof}

\section{Proof of the short canonical-yield lemma}
\label{app:linear-short-proof}

\begin{proof}[Proof of Lemma~\ref{lem:linear-short}]
In the linear-spine SSBNF form of
Proposition~\ref{prop:linear-normal}, each binary spine step has the form
$A\to W_aB$ or $A\to BW_a$, with $W_a\Rightarrow a$, and therefore emits
exactly one terminal outside the continuing nonterminal.

If $X$ is a terminal-wrapper symbol then $|\omega(X)|=1$.  Otherwise choose
a shortest terminal derivation from the typed spine symbol $X$ and write its
spine as $X=X_0,X_1,\ldots,X_t$, where $X_t$ uses a terminal rule.  Its
yield has length $t+1$.  No typed spine symbol can repeat: if
$X_i=X_j$ with $i<j$, deleting the segment between the two occurrences gives
a valid derivation with fewer spine steps and hence a shorter yield.  Thus the
$t+1$ typed spine symbols are distinct non-start symbols, so
$t+1\le N_t$.  Hence $|\omega(X)|\le N_t$ in all cases.
\end{proof}

\begin{thebibliography}{99}

\bibitem{kanazawa1998}
M.~Kanazawa.
\newblock \emph{Learnable Classes of Categorial Grammars}.
\newblock CSLI Publications, 1998.

\bibitem{kanazawa1996}
M.~Kanazawa.
\newblock Identification in the limit of categorial grammars.
\newblock \textit{Journal of Logic, Language and Information},
5(2):115--155, 1996. doi:10.1007/BF00173697.


\bibitem{gold1967}
E.\,M.~Gold.
Language identification in the limit.
\textit{Information and Control}, 10(5):447--474, 1967.

\bibitem{pin1997}
J.-E.~Pin.
\newblock Syntactic semigroups.
\newblock In G.~Rozenberg and A.~Salomaa (eds.),
\textit{Handbook of Formal Languages}, Vol.~1, pp.~679--746,
Springer, 1997.

\bibitem{eilenberg1974}
S.~Eilenberg.
\newblock \emph{Automata, Languages, and Machines}, Vol.~A.
\newblock Academic Press, 1974.

\bibitem{pin2025}
J.-E.~Pin.
\newblock \emph{Mathematical Foundations of Automata Theory}.
\newblock Lecture notes, version of March 24, 2025.
\newblock \url{https://www.irif.fr/~jep/PDF/MPRI/MPRI.pdf}.

\bibitem{higuera1997}
C.~de la Higuera.
Characteristic sets for polynomial grammatical inference.
\textit{Machine Learning}, 27:125--138, 1997.

\bibitem{higuera2010}
C.~de la Higuera.
\textit{Grammatical Inference: Learning Automata and Grammars}.
Cambridge University Press, 2010.


\bibitem{clark-eyraud2007}
A.~Clark and R.~Eyraud.
Polynomial identification in the limit of substitutable
context-free languages.
\textit{Journal of Machine Learning Research},
8:1725--1745, 2007.

\bibitem{wakatsuki-tomita1993}
M.~Wakatsuki and E.~Tomita.
\newblock A fast algorithm for checking the inclusion for very simple
deterministic pushdown automata.
\newblock \textit{IEICE Transactions on Information and Systems},
E76-D(10):1224--1233, 1993.

\bibitem{kuriyama2014}
T.~Kuriyama.
\newblock \emph{Learning Algorithm for Relation-Substitutable Context-Free Languages}.
\newblock arXiv:1409.6247, originally posted in 2014 (v3: November 2014).
\newblock \url{https://arxiv.org/abs/1409.6247}.

\bibitem{yoshinaka2015}
R.~Yoshinaka.
\newblock General perspective on distributionally learnable classes.
\newblock In \emph{Proceedings of the 14th Meeting on the Mathematics
of Language (MoL 14)}, pp.~87--98, 2015.

\bibitem{clark-yoshinaka2016}
A.~Clark and R.~Yoshinaka.
\newblock Distributional learning of context-free and multiple
context-free grammars.
\newblock In \emph{Topics in Grammatical Inference}, Springer, 2016.

\bibitem{coste-garet-nicolas2012}
F.~Coste, G.~Garet, and J.~Nicolas.
\newblock Locally substitutable languages for enhanced inductive leaps.
\newblock In \emph{Proceedings of the Eleventh International Conference
on Grammatical Inference}, PMLR 21:97--111, 2012.

\bibitem{coste-nicolas2019}
F.~Coste and J.~Nicolas.
\newblock Learning local substitutable context-free languages from
positive examples in polynomial time and data by reduction.
\newblock In \emph{Proceedings of the 14th International Conference
on Grammatical Inference}, PMLR 93:155--168, 2019.

\bibitem{yoshinaka2008}
R.~Yoshinaka.
Identification in the limit of $k, l$-substitutable
context-free languages.
In \textit{Grammatical Inference: Algorithms and Applications -- ICGI 2008},
LNCS 5278, pp.~266--279, Springer, 2008.
doi:10.1007/978-3-540-88009-7\_21.


\bibitem{takada1988}
Y.~Takada.
\newblock Grammatical inference for even linear languages based on control sets.
\newblock \textit{Information Processing Letters}, 28(4):193--199, 1988.
doi:10.1016/0020-0190(88)90208-6.

\bibitem{takada1995}
Y.~Takada.
\newblock Learning formal languages based on control sets.
\newblock In K.~P. Jantke and S.~Lange (Eds.), \textit{Algorithmic Learning
for Knowledge-Based Systems}, LNAI 961, pp.~316--339, Springer, 1995.
doi:10.1007/3-540-60217-8\_15.



\bibitem{ginsburg-spanier1966}
S.~Ginsburg and E.~H.~Spanier.
\newblock Finite-turn pushdown automata.
\newblock \textit{SIAM Journal on Control}, 4(3):429--453, 1966.
doi:10.1137/0304034.

\bibitem{coste-typing2004}
F.~Coste, D.~Fredouille, C.~Kermorvant, and C.~de~la~Higuera.
\newblock Introducing domain and typing bias in automata inference.
\newblock In \textit{ICGI 2004}, LNCS 3264, pp.~115--126, Springer, 2004.
doi:10.1007/978-3-540-30195-0\_11.

\bibitem{clark2013}
A.~Clark.
\newblock Learning trees from strings: A strong learning algorithm for some context-free grammars.
\newblock \textit{Journal of Machine Learning Research}, 14:3537--3559, 2013.




\bibitem{eyraud-heinz-yoshinaka2016}
R.~Eyraud, J.~Heinz, and R.~Yoshinaka.
\newblock Efficiency in the identification in the limit learning paradigm.
\newblock In J.~Heinz and J.~M.~Sempere (Eds.),
\emph{Topics in Grammatical Inference}, Springer, 2016.
doi:10.1007/978-3-662-48395-4\_2.

\bibitem{autebert1997}
J.-M.~Autebert, J.~Berstel, and L.~Boasson.
\newblock Context-free languages and pushdown automata.
\newblock In G.~Rozenberg and A.~Salomaa (Eds.),
\emph{Handbook of Formal Languages}, Vol.~1, pp.~111--174,
Springer, 1997.

\bibitem{ginsburg-greibach1966}
S.~Ginsburg and S.~A.~Greibach.
\newblock Deterministic context-free languages.
\newblock \emph{Information and Control}, 9:620--648, 1966.

\bibitem{clark2010congruential}
A.~Clark.
\newblock Distributional learning of some context-free languages with a
minimally adequate teacher.
\newblock In \emph{Proceedings of ICGI 2010}, LNCS 6339, pp.~24--37,
Springer, 2010.
doi:10.1007/978-3-642-15488-1\_4.



\end{thebibliography}
\end{document}